%%%%%%%%%%%%%%%%%%%%%%%%%%%%%%%%%%%%%%%%%%%%%%%%%%%%%%%%%%%%
% 11.8 SYMBOLIC COMPUTATION
% The Composition of Advanced Mathematics
%%%%%%%%%%%%%%%%%%%%%%%%%%%%%%%%%%%%%%%%%%%%%%%%%%%%%%%%%%%%

\section{Symbolic Computation}
\label{sec:symbolic-computation}

\prerequisite{
Algebra, functions, calculus, linear algebra, polynomial arithmetic,
elementary number theory, logic, algorithms, data structures, and basic
programming concepts.
}

\learningobjectives{
By the end of this section, the reader should be able to:
\begin{itemize}
    \item explain the purpose of symbolic computation;
    \item distinguish symbolic and numerical computation;
    \item understand symbolic expressions as structured mathematical
          objects;
    \item recognize expression trees;
    \item distinguish syntactic and mathematical equality;
    \item understand canonical and normal forms;
    \item simplify algebraic expressions symbolically;
    \item understand expansion, collection, and factorization;
    \item perform exact polynomial arithmetic;
    \item understand polynomial greatest common divisors;
    \item recognize the Euclidean algorithm for polynomials;
    \item understand square-free factorization conceptually;
    \item recognize symbolic polynomial factorization;
    \item understand rational-function simplification;
    \item perform partial-fraction decomposition conceptually;
    \item understand exact arithmetic with integers and rational numbers;
    \item recognize algebraic-number representation conceptually;
    \item understand arbitrary-precision symbolic arithmetic;
    \item perform symbolic differentiation;
    \item derive differentiation rules recursively;
    \item understand symbolic integration conceptually;
    \item distinguish elementary antiderivatives from nonelementary
          integrals;
    \item understand definite integration symbolically;
    \item recognize substitution and integration-by-parts strategies;
    \item understand simplification after differentiation and integration;
    \item solve algebraic equations symbolically when possible;
    \item distinguish exact roots from numerical approximations;
    \item understand radicals and polynomial solvability limitations;
    \item recognize resultants and elimination conceptually;
    \item understand substitution-based elimination;
    \item recognize systems of polynomial equations;
    \item understand Gröbner-basis ideas conceptually;
    \item manipulate symbolic matrices;
    \item compute determinants symbolically;
    \item understand exact Gaussian elimination;
    \item recognize symbolic eigenvalue problems;
    \item understand symbolic series expansions;
    \item compute Taylor and Laurent series conceptually;
    \item recognize asymptotic expansions;
    \item understand symbolic limits;
    \item recognize piecewise-defined symbolic expressions;
    \item understand assumptions and domains;
    \item distinguish real and complex simplification;
    \item understand branch cuts conceptually;
    \item recognize pattern matching and rewrite rules;
    \item understand term rewriting;
    \item recognize recursive simplification;
    \item understand expression swell;
    \item recognize common subexpression elimination;
    \item understand exact versus approximate constants;
    \item recognize symbolic manipulation of differential equations;
    \item understand symbolic transforms conceptually;
    \item recognize symbolic summation and products;
    \item understand recurrence relations symbolically;
    \item distinguish symbolic computation from computer algebra systems;
    \item connect symbolic computation with algebra, calculus, number
          theory, differential equations, optimization, theorem proving,
          and scientific computing.
\end{itemize}
}

%%%%%%%%%%%%%%%%%%%%%%%%%%%%%%%%%%%%%%%%%%%%%%%%%%%%%%%%%%%%
\subsection{What Is Symbolic Computation?}
\label{subsec:what-is-symbolic-computation}
%%%%%%%%%%%%%%%%%%%%%%%%%%%%%%%%%%%%%%%%%%%%%%%%%%%%%%%%%%%%

Symbolic computation manipulates mathematical expressions according to
algebraic, logical, and analytic rules.

Instead of immediately approximating quantities numerically, symbolic
methods attempt to preserve exact mathematical structure.

For example,

\[
\boxed{
\sqrt{2}
}
\]

may be stored symbolically rather than approximated by

\[
1.41421356\ldots
\]

Similarly,

\[
\boxed{
\int x^2\,dx
=
\frac{x^3}{3}+C
}
\]

is a symbolic result.

%%%%%%%%%%%%%%%%%%%%%%%%%%%%%%%%%%%%%%%%%%%%%%%%%%%%%%%%%%%%
\subsection{Symbolic Versus Numerical Computation}
\label{subsec:symbolic-versus-numerical}
%%%%%%%%%%%%%%%%%%%%%%%%%%%%%%%%%%%%%%%%%%%%%%%%%%%%%%%%%%%%

Numerical computation might evaluate

\[
\sqrt{2}
\]

as

\[
\boxed{
1.41421356237.
}
\]

Symbolic computation retains

\[
\boxed{
\sqrt{2}.
}
\]

Likewise, a numerical solver may approximate the roots of

\[
x^2-2=0
\]

as

\[
\boxed{
x\approx\pm1.41421356,
}
\]

while a symbolic solver returns

\[
\boxed{
x=\pm\sqrt{2}.
}
\]

%%%%%%%%%%%%%%%%%%%%%%%%%%%%%%%%%%%%%%%%%%%%%%%%%%%%%%%%%%%%
\subsection{Why Symbolic Computation Matters}
\label{subsec:why-symbolic-computation-matters}
%%%%%%%%%%%%%%%%%%%%%%%%%%%%%%%%%%%%%%%%%%%%%%%%%%%%%%%%%%%%

Exact symbolic expressions can reveal mathematical structure that decimal
approximations hide.

Examples include:

\[
\boxed{
\text{factorization},
}
\]

\[
\boxed{
\text{symmetry},
}
\]

\[
\boxed{
\text{cancellation},
}
\]

\[
\boxed{
\text{parameter dependence},
}
\]

and

\[
\boxed{
\text{exact identities}.
}
\]

Symbolic computation is therefore useful both for solving problems and for
understanding them.

%%%%%%%%%%%%%%%%%%%%%%%%%%%%%%%%%%%%%%%%%%%%%%%%%%%%%%%%%%%%
\subsection{Expressions as Mathematical Objects}
\label{subsec:symbolic-expressions}
%%%%%%%%%%%%%%%%%%%%%%%%%%%%%%%%%%%%%%%%%%%%%%%%%%%%%%%%%%%%

A symbolic expression is not merely a string of characters.

It has mathematical structure.

For example,

\[
\boxed{
(x+1)^2
}
\]

contains:

\[
\boxed{
\text{power},
}
\]

\[
\boxed{
\text{sum},
}
\]

\[
\boxed{
\text{symbol},
}
\]

and

\[
\boxed{
\text{constant}.
}
\]

A symbolic system stores and manipulates this structure.

%%%%%%%%%%%%%%%%%%%%%%%%%%%%%%%%%%%%%%%%%%%%%%%%%%%%%%%%%%%%
\subsection{Expression Trees}
\label{subsec:expression-trees}
%%%%%%%%%%%%%%%%%%%%%%%%%%%%%%%%%%%%%%%%%%%%%%%%%%%%%%%%%%%%

The expression

\[
(x+1)^2
\]

can be viewed as a tree whose root is exponentiation.

Conceptually,

\[
\boxed{
\operatorname{Power}
\left(
\operatorname{Add}(x,1),
2
\right).
}
\]

Similarly,

\[
x^2+3x+1
\]

can be represented as a sum of structured terms.

Tree representations allow algorithms to recursively inspect and transform
expressions.

%%%%%%%%%%%%%%%%%%%%%%%%%%%%%%%%%%%%%%%%%%%%%%%%%%%%%%%%%%%%
\subsection{Recursive Expression Processing}
\label{subsec:recursive-expression-processing}
%%%%%%%%%%%%%%%%%%%%%%%%%%%%%%%%%%%%%%%%%%%%%%%%%%%%%%%%%%%%

Many symbolic algorithms work recursively.

For example, to differentiate

\[
f(x)+g(x),
\]

differentiate each child expression:

\[
\boxed{
\frac{d}{dx}
[
f(x)+g(x)
]
=
f'(x)+g'(x).
}
\]

For a product,

\[
\boxed{
\frac{d}{dx}
[
f(x)g(x)
]
=
f'(x)g(x)+f(x)g'(x).
}
\]

Thus symbolic differentiation naturally follows the expression tree.

%%%%%%%%%%%%%%%%%%%%%%%%%%%%%%%%%%%%%%%%%%%%%%%%%%%%%%%%%%%%
\subsection{Syntactic and Mathematical Equality}
\label{subsec:syntactic-mathematical-equality}
%%%%%%%%%%%%%%%%%%%%%%%%%%%%%%%%%%%%%%%%%%%%%%%%%%%%%%%%%%%%

Two expressions can look different but represent the same mathematical
function.

For example,

\[
\boxed{
(x+1)^2
}
\]

and

\[
\boxed{
x^2+2x+1
}
\]

are syntactically different but algebraically equivalent.

Likewise,

\[
\boxed{
\frac{x^2-1}{x-1}
}
\]

simplifies to

\[
x+1
\]

only when

\[
\boxed{
x\neq1.
}
\]

Thus symbolic equality may depend on domains and assumptions.

%%%%%%%%%%%%%%%%%%%%%%%%%%%%%%%%%%%%%%%%%%%%%%%%%%%%%%%%%%%%
\subsection{Canonical Forms}
\label{subsec:canonical-forms}
%%%%%%%%%%%%%%%%%%%%%%%%%%%%%%%%%%%%%%%%%%%%%%%%%%%%%%%%%%%%

A canonical form attempts to represent mathematically equivalent objects
using the same standardized representation.

For example, a polynomial may be written as

\[
\boxed{
a_nx^n
+
a_{n-1}x^{n-1}
+
\cdots
+
a_0.
}
\]

Canonical representations make equality testing easier.

If two expressions reduce to the same canonical form, their equality may
be immediately recognized.

%%%%%%%%%%%%%%%%%%%%%%%%%%%%%%%%%%%%%%%%%%%%%%%%%%%%%%%%%%%%
\subsection{Normal Forms}
\label{subsec:symbolic-normal-forms}
%%%%%%%%%%%%%%%%%%%%%%%%%%%%%%%%%%%%%%%%%%%%%%%%%%%%%%%%%%%%

A normal form is a standardized endpoint of a transformation process.

Different symbolic tasks use different normal forms.

Examples include:

\[
\boxed{
\text{expanded polynomial form},
}
\]

\[
\boxed{
\text{factored form},
}
\]

\[
\boxed{
\text{reduced rational form},
}
\]

and

\[
\boxed{
\text{Gröbner normal form}.
}
\]

There is no single universally best form for every computation.

%%%%%%%%%%%%%%%%%%%%%%%%%%%%%%%%%%%%%%%%%%%%%%%%%%%%%%%%%%%%
\subsection{Simplification}
\label{subsec:symbolic-simplification}
%%%%%%%%%%%%%%%%%%%%%%%%%%%%%%%%%%%%%%%%%%%%%%%%%%%%%%%%%%%%

Symbolic simplification attempts to replace an expression by an equivalent
expression that is easier to interpret or compute.

Examples include:

\[
x+0
\longrightarrow
x,
\]

\[
x\cdot1
\longrightarrow
x,
\]

and

\[
\boxed{
x+x
\longrightarrow
2x.
}
\]

But the notion of ``simpler'' is context dependent.

%%%%%%%%%%%%%%%%%%%%%%%%%%%%%%%%%%%%%%%%%%%%%%%%%%%%%%%%%%%%
\subsection{Ambiguity of Simplification}
\label{subsec:simplification-ambiguity}
%%%%%%%%%%%%%%%%%%%%%%%%%%%%%%%%%%%%%%%%%%%%%%%%%%%%%%%%%%%%

Compare

\[
\boxed{
(x+1)^5
}
\]

with its expanded form.

The factored power is shorter.

The expanded polynomial may be more useful for coefficient extraction.

Similarly,

\[
\boxed{
x^2-1
}
\]

may be simpler for expansion-based operations, while

\[
\boxed{
(x-1)(x+1)
}
\]

reveals roots immediately.

Thus symbolic simplification should depend on the intended task.

%%%%%%%%%%%%%%%%%%%%%%%%%%%%%%%%%%%%%%%%%%%%%%%%%%%%%%%%%%%%
\subsection{Expansion}
\label{subsec:symbolic-expansion}
%%%%%%%%%%%%%%%%%%%%%%%%%%%%%%%%%%%%%%%%%%%%%%%%%%%%%%%%%%%%

Expansion distributes products and powers.

For example,

\[
(x+2)(x-3)
\]

expands to

\[
\boxed{
x^2-x-6.
}
\]

Likewise,

\[
(x+y)^3
\]

expands to

\[
\boxed{
x^3
+
3x^2y
+
3xy^2
+
y^3.
}
\]

%%%%%%%%%%%%%%%%%%%%%%%%%%%%%%%%%%%%%%%%%%%%%%%%%%%%%%%%%%%%
\subsection{Collection of Terms}
\label{subsec:collecting-symbolic-terms}
%%%%%%%%%%%%%%%%%%%%%%%%%%%%%%%%%%%%%%%%%%%%%%%%%%%%%%%%%%%%

Terms may be grouped according to powers of a selected variable.

For example,

\[
ax^2+bx+cx^2+d
\]

can be collected as

\[
\boxed{
(a+c)x^2+bx+d.
}
\]

Collection is especially useful in polynomial and perturbation
calculations.

%%%%%%%%%%%%%%%%%%%%%%%%%%%%%%%%%%%%%%%%%%%%%%%%%%%%%%%%%%%%
\subsection{Factorization}
\label{subsec:symbolic-factorization}
%%%%%%%%%%%%%%%%%%%%%%%%%%%%%%%%%%%%%%%%%%%%%%%%%%%%%%%%%%%%

Factorization reverses expansion.

For example,

\[
x^2-5x+6
\]

factors as

\[
\boxed{
(x-2)(x-3).
}
\]

Factorization can reveal:

\[
\boxed{
\text{roots},
}
\]

\[
\boxed{
\text{multiplicities},
}
\]

and

\[
\boxed{
\text{algebraic structure}.
}
\]

%%%%%%%%%%%%%%%%%%%%%%%%%%%%%%%%%%%%%%%%%%%%%%%%%%%%%%%%%%%%
\subsection{Worked Example: Equivalent Forms}
\label{subsec:worked-symbolic-equivalent-forms}
%%%%%%%%%%%%%%%%%%%%%%%%%%%%%%%%%%%%%%%%%%%%%%%%%%%%%%%%%%%%

\begin{workedexamplebox}{Expand and Refactor}

Consider

\[
(x-2)(x+5).
\]

Expanding gives

\[
x^2+5x-2x-10.
\]

Therefore,

\[
x^2+3x-10.
\]

Factoring the result gives

\[
(x-2)(x+5).
\]

\begin{answerbox}
\[
\boxed{
(x-2)(x+5)
=
x^2+3x-10.
}
\]
\end{answerbox}

The two forms are mathematically equivalent but useful for different
purposes.

\end{workedexamplebox}

%%%%%%%%%%%%%%%%%%%%%%%%%%%%%%%%%%%%%%%%%%%%%%%%%%%%%%%%%%%%
\subsection{Exact Integer Arithmetic}
\label{subsec:exact-integer-arithmetic}
%%%%%%%%%%%%%%%%%%%%%%%%%%%%%%%%%%%%%%%%%%%%%%%%%%%%%%%%%%%%

Integers can be represented exactly.

Unlike fixed-width machine integers, symbolic integer arithmetic can use
arbitrarily many digits, limited primarily by computational resources.

Thus symbolic systems can compute quantities such as

\[
\boxed{
100!
}
\]

exactly.

No floating-point approximation is required.

%%%%%%%%%%%%%%%%%%%%%%%%%%%%%%%%%%%%%%%%%%%%%%%%%%%%%%%%%%%%
\subsection{Exact Rational Arithmetic}
\label{subsec:exact-rational-arithmetic}
%%%%%%%%%%%%%%%%%%%%%%%%%%%%%%%%%%%%%%%%%%%%%%%%%%%%%%%%%%%%

A rational number is represented as

\[
\boxed{
\frac{p}{q},
\qquad
p,q\in\mathbb{Z},
\quad
q\neq0.
}
\]

The fraction is often normalized so that

\[
\gcd(p,q)=1.
\]

For example,

\[
\frac{8}{12}
\]

is reduced to

\[
\boxed{
\frac23.
}
\]

%%%%%%%%%%%%%%%%%%%%%%%%%%%%%%%%%%%%%%%%%%%%%%%%%%%%%%%%%%%%
\subsection{Exact Versus Floating-Point Constants}
\label{subsec:exact-vs-floating-constants}
%%%%%%%%%%%%%%%%%%%%%%%%%%%%%%%%%%%%%%%%%%%%%%%%%%%%%%%%%%%%

The expression

\[
\boxed{
\frac13
}
\]

is exact.

The decimal

\[
\boxed{
0.3333333333
}
\]

is approximate.

Likewise,

\[
\boxed{
\pi
}
\]

is an exact symbolic constant, while

\[
\boxed{
3.14159
}
\]

is a finite numerical approximation.

Mixing exact and approximate inputs may cause a symbolic system to switch
to approximate arithmetic.

%%%%%%%%%%%%%%%%%%%%%%%%%%%%%%%%%%%%%%%%%%%%%%%%%%%%%%%%%%%%
\subsection{Polynomial Representation}
\label{subsec:symbolic-polynomial-representation}
%%%%%%%%%%%%%%%%%%%%%%%%%%%%%%%%%%%%%%%%%%%%%%%%%%%%%%%%%%%%

A polynomial

\[
p(x)
=
a_0+a_1x+\cdots+a_nx^n
\]

may be represented by its coefficients:

\[
\boxed{
(a_0,a_1,\ldots,a_n).
}
\]

Alternatively, a sparse representation stores only nonzero terms:

\[
\boxed{
\{
(k,a_k)
:
a_k\neq0
\}.
}
\]

The best representation depends on polynomial density and degree.

%%%%%%%%%%%%%%%%%%%%%%%%%%%%%%%%%%%%%%%%%%%%%%%%%%%%%%%%%%%%
\subsection{Polynomial Addition}
\label{subsec:symbolic-polynomial-addition}
%%%%%%%%%%%%%%%%%%%%%%%%%%%%%%%%%%%%%%%%%%%%%%%%%%%%%%%%%%%%

If

\[
p(x)=\sum_{k=0}^{m}a_kx^k
\]

and

\[
q(x)=\sum_{k=0}^{n}b_kx^k,
\]

then

\[
\boxed{
p(x)+q(x)
=
\sum_k
(a_k+b_k)x^k.
}
\]

Symbolic polynomial addition therefore reduces to coefficient addition.

%%%%%%%%%%%%%%%%%%%%%%%%%%%%%%%%%%%%%%%%%%%%%%%%%%%%%%%%%%%%
\subsection{Polynomial Multiplication}
\label{subsec:symbolic-polynomial-multiplication}
%%%%%%%%%%%%%%%%%%%%%%%%%%%%%%%%%%%%%%%%%%%%%%%%%%%%%%%%%%%%

If

\[
p(x)
=
\sum_i
a_ix^i
\]

and

\[
q(x)
=
\sum_j
b_jx^j,
\]

then

\[
\boxed{
p(x)q(x)
=
\sum_k
\left(
\sum_{i+j=k}
a_ib_j
\right)
x^k.
}
\]

The coefficients are obtained by discrete convolution.

%%%%%%%%%%%%%%%%%%%%%%%%%%%%%%%%%%%%%%%%%%%%%%%%%%%%%%%%%%%%
\subsection{Polynomial Division}
\label{subsec:symbolic-polynomial-division}
%%%%%%%%%%%%%%%%%%%%%%%%%%%%%%%%%%%%%%%%%%%%%%%%%%%%%%%%%%%%

For polynomials \(f\) and nonzero \(g\) over a field, there exist unique
polynomials \(q\) and \(r\) such that

\[
\boxed{
f
=
qg+r,
}
\]

with

\[
\boxed{
\deg r
<
\deg g.
}
\]

This is the polynomial division algorithm.

%%%%%%%%%%%%%%%%%%%%%%%%%%%%%%%%%%%%%%%%%%%%%%%%%%%%%%%%%%%%
\subsection{Polynomial Greatest Common Divisor}
\label{subsec:polynomial-gcd}
%%%%%%%%%%%%%%%%%%%%%%%%%%%%%%%%%%%%%%%%%%%%%%%%%%%%%%%%%%%%

A greatest common divisor of polynomials \(f\) and \(g\) is a polynomial
\(d\) satisfying

\[
d\mid f,
\]

\[
d\mid g,
\]

and every common divisor divides \(d\).

Over a field, the gcd is usually normalized to be monic.

%%%%%%%%%%%%%%%%%%%%%%%%%%%%%%%%%%%%%%%%%%%%%%%%%%%%%%%%%%%%
\subsection{Polynomial Euclidean Algorithm}
\label{subsec:polynomial-euclidean-algorithm}
%%%%%%%%%%%%%%%%%%%%%%%%%%%%%%%%%%%%%%%%%%%%%%%%%%%%%%%%%%%%

The Euclidean algorithm repeatedly computes remainders:

\[
f=q_0g+r_1,
\]

\[
g=q_1r_1+r_2,
\]

and so on.

When the remainder becomes zero, the final nonzero remainder is the gcd,
up to multiplication by a nonzero constant.

%%%%%%%%%%%%%%%%%%%%%%%%%%%%%%%%%%%%%%%%%%%%%%%%%%%%%%%%%%%%
\subsection{Worked Example: Polynomial GCD}
\label{subsec:worked-polynomial-gcd}
%%%%%%%%%%%%%%%%%%%%%%%%%%%%%%%%%%%%%%%%%%%%%%%%%%%%%%%%%%%%

\begin{workedexamplebox}{Finding a Common Polynomial Factor}

Consider

\[
f(x)
=
x^3-x
\]

and

\[
g(x)
=
x^2-1.
\]

Factor:

\[
f(x)
=
x(x-1)(x+1),
\]

and

\[
g(x)
=
(x-1)(x+1).
\]

Therefore,

\begin{answerbox}
\[
\boxed{
\gcd(f,g)
=
x^2-1
}
\]
\end{answerbox}

when the gcd is chosen monic.

\end{workedexamplebox}

%%%%%%%%%%%%%%%%%%%%%%%%%%%%%%%%%%%%%%%%%%%%%%%%%%%%%%%%%%%%
\subsection{Square-Free Polynomials}
\label{subsec:square-free-polynomials}
%%%%%%%%%%%%%%%%%%%%%%%%%%%%%%%%%%%%%%%%%%%%%%%%%%%%%%%%%%%%

A polynomial is square-free if it has no repeated irreducible factor.

Repeated factors can be detected through

\[
\boxed{
\gcd
(
f,f'
).
}
\]

If

\[
\gcd(f,f')=1,
\]

then \(f\) is square-free over a characteristic-zero field.

%%%%%%%%%%%%%%%%%%%%%%%%%%%%%%%%%%%%%%%%%%%%%%%%%%%%%%%%%%%%
\subsection{Repeated Roots}
\label{subsec:symbolic-repeated-roots}
%%%%%%%%%%%%%%%%%%%%%%%%%%%%%%%%%%%%%%%%%%%%%%%%%%%%%%%%%%%%

Suppose

\[
f(x)
=
(x-a)^m g(x),
\]

where

\[
g(a)\neq0.
\]

Then

\[
a
\]

is a root of multiplicity \(m\).

A repeated root is simultaneously a root of

\[
f
\]

and

\[
f'.
\]

Thus polynomial gcd computations detect root multiplicity algebraically.

%%%%%%%%%%%%%%%%%%%%%%%%%%%%%%%%%%%%%%%%%%%%%%%%%%%%%%%%%%%%
\subsection{Polynomial Factorization}
\label{subsec:polynomial-factorization-symbolic}
%%%%%%%%%%%%%%%%%%%%%%%%%%%%%%%%%%%%%%%%%%%%%%%%%%%%%%%%%%%%

Symbolic polynomial factorization depends strongly on the coefficient
domain.

For example,

\[
x^2-2
\]

is irreducible over

\[
\mathbb{Q},
\]

but factors over

\[
\mathbb{R}
\]

as

\[
\boxed{
(x-\sqrt{2})(x+\sqrt{2}).
}
\]

Thus factorization must specify the underlying number system.

%%%%%%%%%%%%%%%%%%%%%%%%%%%%%%%%%%%%%%%%%%%%%%%%%%%%%%%%%%%%
\subsection{Rational Functions}
\label{subsec:symbolic-rational-functions}
%%%%%%%%%%%%%%%%%%%%%%%%%%%%%%%%%%%%%%%%%%%%%%%%%%%%%%%%%%%%

A rational function has form

\[
\boxed{
R(x)
=
\frac{
p(x)
}{
q(x)
},
}
\]

where \(p\) and \(q\) are polynomials and

\[
q(x)\neq0.
\]

Common polynomial factors can be canceled algebraically, but excluded
domain points must be remembered.

%%%%%%%%%%%%%%%%%%%%%%%%%%%%%%%%%%%%%%%%%%%%%%%%%%%%%%%%%%%%
\subsection{Rational Simplification}
\label{subsec:rational-simplification}
%%%%%%%%%%%%%%%%%%%%%%%%%%%%%%%%%%%%%%%%%%%%%%%%%%%%%%%%%%%%

Consider

\[
\frac{
x^2-1
}{
x^2-2x+1
}.
\]

Factor:

\[
x^2-1
=
(x-1)(x+1),
\]

and

\[
x^2-2x+1
=
(x-1)^2.
\]

Thus,

\[
\boxed{
\frac{
x^2-1
}{
x^2-2x+1
}
=
\frac{
x+1
}{
x-1
},
\qquad
x\neq1.
}
\]

Domain restrictions remain part of the original expression.

%%%%%%%%%%%%%%%%%%%%%%%%%%%%%%%%%%%%%%%%%%%%%%%%%%%%%%%%%%%%
\subsection{Partial Fractions}
\label{subsec:symbolic-partial-fractions}
%%%%%%%%%%%%%%%%%%%%%%%%%%%%%%%%%%%%%%%%%%%%%%%%%%%%%%%%%%%%

A rational function may be decomposed into simpler rational terms.

For example,

\[
\frac{
1
}{
(x-1)(x+2)
}
\]

can be written

\[
\boxed{
\frac{A}{x-1}
+
\frac{B}{x+2}.
}
\]

Solving for the constants gives a representation useful for integration,
inverse transforms, and differential equations.

%%%%%%%%%%%%%%%%%%%%%%%%%%%%%%%%%%%%%%%%%%%%%%%%%%%%%%%%%%%%
\subsection{Worked Example: Partial Fractions}
\label{subsec:worked-partial-fractions}
%%%%%%%%%%%%%%%%%%%%%%%%%%%%%%%%%%%%%%%%%%%%%%%%%%%%%%%%%%%%

\begin{workedexamplebox}{Decomposing a Rational Function}

Write

\[
\frac1{(x-1)(x+2)}
=
\frac{A}{x-1}
+
\frac{B}{x+2}.
\]

Multiply by

\[
(x-1)(x+2):
\]

\[
1
=
A(x+2)
+
B(x-1).
\]

Set

\[
x=1:
\]

\[
1=3A,
\]

so

\[
A=\frac13.
\]

Set

\[
x=-2:
\]

\[
1=-3B,
\]

so

\[
B=-\frac13.
\]

Therefore,

\begin{answerbox}
\[
\boxed{
\frac1{(x-1)(x+2)}
=
\frac1{3(x-1)}
-
\frac1{3(x+2)}.
}
\]
\end{answerbox}

\end{workedexamplebox}

%%%%%%%%%%%%%%%%%%%%%%%%%%%%%%%%%%%%%%%%%%%%%%%%%%%%%%%%%%%%
\subsection{Symbolic Differentiation}
\label{subsec:symbolic-differentiation}
%%%%%%%%%%%%%%%%%%%%%%%%%%%%%%%%%%%%%%%%%%%%%%%%%%%%%%%%%%%%

Symbolic differentiation applies differentiation rules recursively.

For constants,

\[
\boxed{
\frac{d}{dx}c
=
0.
}
\]

For the variable,

\[
\boxed{
\frac{d}{dx}x
=
1.
}
\]

For powers,

\[
\boxed{
\frac{d}{dx}x^n
=
nx^{n-1}.
}
\]

%%%%%%%%%%%%%%%%%%%%%%%%%%%%%%%%%%%%%%%%%%%%%%%%%%%%%%%%%%%%
\subsection{Recursive Differentiation Rules}
\label{subsec:recursive-differentiation-rules}
%%%%%%%%%%%%%%%%%%%%%%%%%%%%%%%%%%%%%%%%%%%%%%%%%%%%%%%%%%%%

For a sum,

\[
\boxed{
D(f+g)
=
D(f)+D(g).
}
\]

For a product,

\[
\boxed{
D(fg)
=
D(f)g
+
fD(g).
}
\]

For a quotient,

\[
\boxed{
D
\left(
\frac{f}{g}
\right)
=
\frac{
D(f)g-fD(g)
}{
g^2
}.
}
\]

For composition,

\[
\boxed{
D
(
f(g(x))
)
=
f'(g(x))g'(x).
}
\]

%%%%%%%%%%%%%%%%%%%%%%%%%%%%%%%%%%%%%%%%%%%%%%%%%%%%%%%%%%%%
\subsection{Worked Example: Symbolic Differentiation}
\label{subsec:worked-symbolic-differentiation}
%%%%%%%%%%%%%%%%%%%%%%%%%%%%%%%%%%%%%%%%%%%%%%%%%%%%%%%%%%%%

\begin{workedexamplebox}{Differentiating a Composite Product}

Let

\[
f(x)
=
x^2\sin(x^3).
\]

Using the product rule,

\[
f'(x)
=
2x\sin(x^3)
+
x^2
\frac{d}{dx}
\sin(x^3).
\]

Using the chain rule,

\[
\frac{d}{dx}
\sin(x^3)
=
3x^2\cos(x^3).
\]

Therefore,

\begin{answerbox}
\[
\boxed{
f'(x)
=
2x\sin(x^3)
+
3x^4\cos(x^3).
}
\]
\end{answerbox}

\end{workedexamplebox}

%%%%%%%%%%%%%%%%%%%%%%%%%%%%%%%%%%%%%%%%%%%%%%%%%%%%%%%%%%%%
\subsection{Derivative Simplification}
\label{subsec:derivative-simplification}
%%%%%%%%%%%%%%%%%%%%%%%%%%%%%%%%%%%%%%%%%%%%%%%%%%%%%%%%%%%%

Direct recursive differentiation can produce unnecessarily large
expressions.

For example, applying product and chain rules may create repeated factors.

A symbolic system therefore often follows differentiation with:

\[
\boxed{
\text{collection},
}
\]

\[
\boxed{
\text{factorization},
}
\]

or

\[
\boxed{
\text{common-subexpression reduction}.
}
\]

%%%%%%%%%%%%%%%%%%%%%%%%%%%%%%%%%%%%%%%%%%%%%%%%%%%%%%%%%%%%
\subsection{Higher Derivatives}
\label{subsec:symbolic-higher-derivatives}
%%%%%%%%%%%%%%%%%%%%%%%%%%%%%%%%%%%%%%%%%%%%%%%%%%%%%%%%%%%%

Higher derivatives can be obtained recursively:

\[
\boxed{
f^{(n)}
=
D^nf.
}
\]

For example,

\[
f(x)=e^{ax}
\]

satisfies

\[
\boxed{
f^{(n)}(x)
=
a^ne^{ax}.
}
\]

Symbolic computation can detect and exploit such recurring patterns.

%%%%%%%%%%%%%%%%%%%%%%%%%%%%%%%%%%%%%%%%%%%%%%%%%%%%%%%%%%%%
\subsection{Symbolic Integration}
\label{subsec:symbolic-integration}
%%%%%%%%%%%%%%%%%%%%%%%%%%%%%%%%%%%%%%%%%%%%%%%%%%%%%%%%%%%%

Symbolic integration seeks an expression \(F\) satisfying

\[
\boxed{
F'(x)=f(x).
}
\]

Then

\[
\boxed{
\int f(x)\,dx
=
F(x)+C.
}
\]

Unlike differentiation, there is no simple universal local recursion that
solves every elementary integration problem.

%%%%%%%%%%%%%%%%%%%%%%%%%%%%%%%%%%%%%%%%%%%%%%%%%%%%%%%%%%%%
\subsection{Basic Integration Rules}
\label{subsec:symbolic-integration-rules}
%%%%%%%%%%%%%%%%%%%%%%%%%%%%%%%%%%%%%%%%%%%%%%%%%%%%%%%%%%%%

Examples include

\[
\boxed{
\int x^n\,dx
=
\frac{
x^{n+1}
}{
n+1
}
+
C,
\qquad
n\neq-1,
}
\]

\[
\boxed{
\int e^x\,dx
=
e^x+C,
}
\]

and

\[
\boxed{
\int \frac1x\,dx
=
\ln|x|+C
}
\]

over real intervals not containing \(0\).

%%%%%%%%%%%%%%%%%%%%%%%%%%%%%%%%%%%%%%%%%%%%%%%%%%%%%%%%%%%%
\subsection{Substitution}
\label{subsec:symbolic-substitution-integration}
%%%%%%%%%%%%%%%%%%%%%%%%%%%%%%%%%%%%%%%%%%%%%%%%%%%%%%%%%%%%

If the integrand contains

\[
f(g(x))g'(x),
\]

set

\[
u=g(x).
\]

Then

\[
du=g'(x)\,dx.
\]

Thus,

\[
\boxed{
\int
f(g(x))g'(x)\,dx
=
\int
f(u)\,du.
}
\]

Symbolic integration systems search for patterns resembling this structure.

%%%%%%%%%%%%%%%%%%%%%%%%%%%%%%%%%%%%%%%%%%%%%%%%%%%%%%%%%%%%
\subsection{Integration by Parts}
\label{subsec:symbolic-integration-by-parts}
%%%%%%%%%%%%%%%%%%%%%%%%%%%%%%%%%%%%%%%%%%%%%%%%%%%%%%%%%%%%

Integration by parts is

\[
\boxed{
\int u\,dv
=
uv
-
\int v\,du.
}
\]

Symbolically, the challenge is selecting \(u\) and \(dv\) in a way that
reduces the complexity of the remaining integral.

This choice requires heuristic or algorithmic strategy.

%%%%%%%%%%%%%%%%%%%%%%%%%%%%%%%%%%%%%%%%%%%%%%%%%%%%%%%%%%%%
\subsection{Worked Example: Symbolic Integration}
\label{subsec:worked-symbolic-integration}
%%%%%%%%%%%%%%%%%%%%%%%%%%%%%%%%%%%%%%%%%%%%%%%%%%%%%%%%%%%%

\begin{workedexamplebox}{Integration by Parts}

Consider

\[
\int xe^x\,dx.
\]

Choose

\[
u=x,
\qquad
dv=e^x\,dx.
\]

Then

\[
du=dx,
\qquad
v=e^x.
\]

Therefore,

\[
\int xe^x\,dx
=
xe^x
-
\int e^x\,dx.
\]

Hence,

\begin{answerbox}
\[
\boxed{
\int xe^x\,dx
=
e^x(x-1)+C.
}
\]
\end{answerbox}

\end{workedexamplebox}

%%%%%%%%%%%%%%%%%%%%%%%%%%%%%%%%%%%%%%%%%%%%%%%%%%%%%%%%%%%%
\subsection{Nonelementary Antiderivatives}
\label{subsec:nonelementary-antiderivatives}
%%%%%%%%%%%%%%%%%%%%%%%%%%%%%%%%%%%%%%%%%%%%%%%%%%%%%%%%%%%%

Not every elementary function has an elementary antiderivative.

For example,

\[
\boxed{
\int e^{-x^2}\,dx
}
\]

cannot be expressed using only a finite combination of elementary
functions.

Instead, one introduces special functions.

This illustrates a fundamental limitation of symbolic integration.

%%%%%%%%%%%%%%%%%%%%%%%%%%%%%%%%%%%%%%%%%%%%%%%%%%%%%%%%%%%%
\subsection{Special Functions}
\label{subsec:symbolic-special-functions}
%%%%%%%%%%%%%%%%%%%%%%%%%%%%%%%%%%%%%%%%%%%%%%%%%%%%%%%%%%%%

Symbolic systems may express results using named special functions.

Examples include:

\[
\boxed{
\operatorname{erf}(x),
}
\]

\[
\boxed{
\Gamma(x),
}
\]

\[
\boxed{
\operatorname{Li}_s(x),
}
\]

and Bessel functions.

Special functions extend the symbolic vocabulary available for exact
solutions.

%%%%%%%%%%%%%%%%%%%%%%%%%%%%%%%%%%%%%%%%%%%%%%%%%%%%%%%%%%%%
\subsection{Definite Symbolic Integration}
\label{subsec:definite-symbolic-integration}
%%%%%%%%%%%%%%%%%%%%%%%%%%%%%%%%%%%%%%%%%%%%%%%%%%%%%%%%%%%%

A definite integral may be evaluated by:

\[
\boxed{
\text{finding an antiderivative},
}
\]

\[
\boxed{
\text{transforming the integral},
}
\]

or

\[
\boxed{
\text{using known identities}.
}
\]

For example,

\[
\boxed{
\int_0^1x^2\,dx
=
\frac13.
}
\]

Some definite integrals have exact values even when no elementary
indefinite antiderivative exists.

%%%%%%%%%%%%%%%%%%%%%%%%%%%%%%%%%%%%%%%%%%%%%%%%%%%%%%%%%%%%
\subsection{Symbolic Limits}
\label{subsec:symbolic-limits}
%%%%%%%%%%%%%%%%%%%%%%%%%%%%%%%%%%%%%%%%%%%%%%%%%%%%%%%%%%%%

Symbolic limit computation seeks

\[
\boxed{
\lim_{x\to a}f(x).
}
\]

Possible methods include:

\[
\boxed{
\text{direct substitution},
}
\]

\[
\boxed{
\text{factor cancellation},
}
\]

\[
\boxed{
\text{series expansion},
}
\]

and

\[
\boxed{
\text{algebraic transformation}.
}
\]

%%%%%%%%%%%%%%%%%%%%%%%%%%%%%%%%%%%%%%%%%%%%%%%%%%%%%%%%%%%%
\subsection{Worked Example: Symbolic Limit}
\label{subsec:worked-symbolic-limit}
%%%%%%%%%%%%%%%%%%%%%%%%%%%%%%%%%%%%%%%%%%%%%%%%%%%%%%%%%%%%

\begin{workedexamplebox}{Removing a Removable Singularity}

Consider

\[
\lim_{x\to1}
\frac{
x^2-1
}{
x-1
}.
\]

Factor:

\[
x^2-1
=
(x-1)(x+1).
\]

For

\[
x\neq1,
\]

\[
\frac{
x^2-1
}{
x-1
}
=
x+1.
\]

Therefore,

\begin{answerbox}
\[
\boxed{
\lim_{x\to1}
\frac{
x^2-1
}{
x-1
}
=
2.
}
\]
\end{answerbox}

\end{workedexamplebox}

%%%%%%%%%%%%%%%%%%%%%%%%%%%%%%%%%%%%%%%%%%%%%%%%%%%%%%%%%%%%
\subsection{Series Expansion}
\label{subsec:symbolic-series-expansion}
%%%%%%%%%%%%%%%%%%%%%%%%%%%%%%%%%%%%%%%%%%%%%%%%%%%%%%%%%%%%

For sufficiently smooth \(f\), a Taylor expansion about \(a\) is

\[
\boxed{
f(x)
=
\sum_{n=0}^{\infty}
\frac{
f^{(n)}(a)
}{
n!
}
(x-a)^n.
}
\]

A symbolic system may produce a truncated series:

\[
\boxed{
f(x)
=
a_0
+
a_1(x-a)
+
\cdots
+
a_N(x-a)^N
+
O((x-a)^{N+1}).
}
\]

%%%%%%%%%%%%%%%%%%%%%%%%%%%%%%%%%%%%%%%%%%%%%%%%%%%%%%%%%%%%
\subsection{Worked Example: Exponential Series}
\label{subsec:worked-exponential-series}
%%%%%%%%%%%%%%%%%%%%%%%%%%%%%%%%%%%%%%%%%%%%%%%%%%%%%%%%%%%%

\begin{workedexamplebox}{Series of \(e^x\)}

Since every derivative of \(e^x\) is \(e^x\),

\[
f^{(n)}(0)=1.
\]

Therefore,

\[
e^x
=
1
+
x
+
\frac{x^2}{2!}
+
\frac{x^3}{3!}
+
\cdots.
\]

Through cubic order,

\begin{answerbox}
\[
\boxed{
e^x
=
1+x+\frac{x^2}{2}+\frac{x^3}{6}+O(x^4).
}
\]
\end{answerbox}

\end{workedexamplebox}

%%%%%%%%%%%%%%%%%%%%%%%%%%%%%%%%%%%%%%%%%%%%%%%%%%%%%%%%%%%%
\subsection{Laurent Series}
\label{subsec:symbolic-laurent-series}
%%%%%%%%%%%%%%%%%%%%%%%%%%%%%%%%%%%%%%%%%%%%%%%%%%%%%%%%%%%%

Near singular points, negative powers may be required.

A Laurent series has form

\[
\boxed{
\sum_{n=-\infty}^{\infty}
a_n
(x-a)^n.
}
\]

For example,

\[
\boxed{
\frac{e^x}{x}
=
\frac1x
+
1
+
\frac{x}{2}
+
\frac{x^2}{6}
+
\cdots.
}
\]

Laurent expansions are important in complex analysis.

%%%%%%%%%%%%%%%%%%%%%%%%%%%%%%%%%%%%%%%%%%%%%%%%%%%%%%%%%%%%
\subsection{Asymptotic Expansion}
\label{subsec:symbolic-asymptotic-expansion}
%%%%%%%%%%%%%%%%%%%%%%%%%%%%%%%%%%%%%%%%%%%%%%%%%%%%%%%%%%%%

An asymptotic expansion approximates behavior in a limiting regime.

For example,

\[
\boxed{
\frac1{1+x}
\sim
1-x+x^2-x^3+\cdots
}
\]

as

\[
x\to0.
\]

An asymptotic series need not converge for fixed nonzero \(x\).

Its meaning concerns successive approximation in the stated limit.

%%%%%%%%%%%%%%%%%%%%%%%%%%%%%%%%%%%%%%%%%%%%%%%%%%%%%%%%%%%%
\subsection{Symbolic Equation Solving}
\label{subsec:symbolic-equation-solving}
%%%%%%%%%%%%%%%%%%%%%%%%%%%%%%%%%%%%%%%%%%%%%%%%%%%%%%%%%%%%

A symbolic equation solver seeks exact expressions satisfying

\[
\boxed{
F(x)=0.
}
\]

For a linear equation

\[
ax+b=0,
\]

with

\[
a\neq0,
\]

the solution is

\[
\boxed{
x=-\frac{b}{a}.
}
\]

%%%%%%%%%%%%%%%%%%%%%%%%%%%%%%%%%%%%%%%%%%%%%%%%%%%%%%%%%%%%
\subsection{Quadratic Formula}
\label{subsec:symbolic-quadratic-formula}
%%%%%%%%%%%%%%%%%%%%%%%%%%%%%%%%%%%%%%%%%%%%%%%%%%%%%%%%%%%%

For

\[
ax^2+bx+c=0,
\qquad
a\neq0,
\]

the symbolic roots are

\[
\boxed{
x
=
\frac{
-b\pm\sqrt{b^2-4ac}
}{
2a
}.
}
\]

The discriminant

\[
\boxed{
\Delta
=
b^2-4ac
}
\]

controls the root structure over the real numbers.

%%%%%%%%%%%%%%%%%%%%%%%%%%%%%%%%%%%%%%%%%%%%%%%%%%%%%%%%%%%%
\subsection{Exact and Numerical Roots}
\label{subsec:exact-numerical-roots}
%%%%%%%%%%%%%%%%%%%%%%%%%%%%%%%%%%%%%%%%%%%%%%%%%%%%%%%%%%%%

For

\[
x^5-x-1=0,
\]

a symbolic system may not produce a root using elementary radicals.

A numerical solver can still approximate roots.

Thus symbolic and numerical methods complement each other:

\[
\boxed{
\text{symbolic where possible}
+
\text{numerical where necessary}.
}
\]

%%%%%%%%%%%%%%%%%%%%%%%%%%%%%%%%%%%%%%%%%%%%%%%%%%%%%%%%%%%%
\subsection{Limits of Radical Solutions}
\label{subsec:limits-of-radical-solutions}
%%%%%%%%%%%%%%%%%%%%%%%%%%%%%%%%%%%%%%%%%%%%%%%%%%%%%%%%%%%%

General polynomial equations of degree five or higher cannot always be
solved using radicals.

This is an algebraic limitation rather than a failure of a particular
software system.

Symbolic systems may instead represent roots implicitly or using special
algebraic objects.

%%%%%%%%%%%%%%%%%%%%%%%%%%%%%%%%%%%%%%%%%%%%%%%%%%%%%%%%%%%%
\subsection{Root Objects}
\label{subsec:symbolic-root-objects}
%%%%%%%%%%%%%%%%%%%%%%%%%%%%%%%%%%%%%%%%%%%%%%%%%%%%%%%%%%%%

If no convenient closed form exists, a symbolic system may represent a
root as:

\[
\boxed{
\text{the }k\text{-th root of }p(x),
}
\]

possibly together with an isolating interval or ordering information.

Such an object can remain exact even without a radical formula.

%%%%%%%%%%%%%%%%%%%%%%%%%%%%%%%%%%%%%%%%%%%%%%%%%%%%%%%%%%%%
\subsection{Systems of Equations}
\label{subsec:symbolic-systems-equations}
%%%%%%%%%%%%%%%%%%%%%%%%%%%%%%%%%%%%%%%%%%%%%%%%%%%%%%%%%%%%

A system may contain several equations:

\[
\boxed{
\begin{aligned}
f_1(x_1,\ldots,x_n)&=0,\\
f_2(x_1,\ldots,x_n)&=0,\\
&\vdots\\
f_m(x_1,\ldots,x_n)&=0.
\end{aligned}
}
\]

Linear systems can be solved exactly by symbolic elimination.

Polynomial nonlinear systems require more advanced algebraic techniques.

%%%%%%%%%%%%%%%%%%%%%%%%%%%%%%%%%%%%%%%%%%%%%%%%%%%%%%%%%%%%
\subsection{Substitution Elimination}
\label{subsec:symbolic-substitution-elimination}
%%%%%%%%%%%%%%%%%%%%%%%%%%%%%%%%%%%%%%%%%%%%%%%%%%%%%%%%%%%%

Consider

\[
x+y=3,
\]

\[
x^2+y^2=5.
\]

From the first equation,

\[
y=3-x.
\]

Substituting into the second gives

\[
x^2+(3-x)^2=5.
\]

Thus a system of two equations is reduced to one polynomial equation.

%%%%%%%%%%%%%%%%%%%%%%%%%%%%%%%%%%%%%%%%%%%%%%%%%%%%%%%%%%%%
\subsection{Worked Example: Symbolic System}
\label{subsec:worked-symbolic-system}
%%%%%%%%%%%%%%%%%%%%%%%%%%%%%%%%%%%%%%%%%%%%%%%%%%%%%%%%%%%%

\begin{workedexamplebox}{Solving by Elimination}

Solve

\[
x+y=3,
\]

\[
x^2+y^2=5.
\]

Use

\[
y=3-x.
\]

Then

\[
x^2+(3-x)^2=5.
\]

Expanding,

\[
2x^2-6x+4=0.
\]

Divide by \(2\):

\[
x^2-3x+2=0.
\]

Therefore,

\[
(x-1)(x-2)=0.
\]

So

\[
x=1
\]

or

\[
x=2.
\]

Correspondingly,

\[
y=2
\]

or

\[
y=1.
\]

\begin{answerbox}
\[
\boxed{
(x,y)
=
(1,2)
\quad
\text{or}
\quad
(2,1).
}
\]
\end{answerbox}

\end{workedexamplebox}

%%%%%%%%%%%%%%%%%%%%%%%%%%%%%%%%%%%%%%%%%%%%%%%%%%%%%%%%%%%%
\subsection{Resultants}
\label{subsec:symbolic-resultants}
%%%%%%%%%%%%%%%%%%%%%%%%%%%%%%%%%%%%%%%%%%%%%%%%%%%%%%%%%%%%

Given two polynomials

\[
f(x,y)
\]

and

\[
g(x,y),
\]

the resultant with respect to \(y\),

\[
\boxed{
\operatorname{Res}_y(f,g),
}
\]

eliminates \(y\).

If \(f\) and \(g\) have a common solution in \(y\), the resultant vanishes
at the corresponding \(x\)-value under appropriate conditions.

Resultants provide a systematic elimination method.

%%%%%%%%%%%%%%%%%%%%%%%%%%%%%%%%%%%%%%%%%%%%%%%%%%%%%%%%%%%%
\subsection{Polynomial Ideals Preview}
\label{subsec:polynomial-ideals-symbolic}
%%%%%%%%%%%%%%%%%%%%%%%%%%%%%%%%%%%%%%%%%%%%%%%%%%%%%%%%%%%%

Given polynomials

\[
f_1,\ldots,f_m,
\]

their polynomial ideal consists of all combinations

\[
\boxed{
h_1f_1
+
\cdots
+
h_mf_m,
}
\]

where the \(h_i\) are polynomials.

Polynomial systems can therefore be studied through the ideal generated by
their equations.

This connects symbolic computation with abstract algebra and algebraic
geometry.

%%%%%%%%%%%%%%%%%%%%%%%%%%%%%%%%%%%%%%%%%%%%%%%%%%%%%%%%%%%%
\subsection{Gröbner Bases Preview}
\label{subsec:groebner-bases-preview}
%%%%%%%%%%%%%%%%%%%%%%%%%%%%%%%%%%%%%%%%%%%%%%%%%%%%%%%%%%%%

A Gröbner basis is a specially organized generating set for a polynomial
ideal relative to a chosen monomial ordering.

It can transform a difficult multivariable polynomial system into a form
better suited for:

\[
\boxed{
\text{elimination},
}
\]

\[
\boxed{
\text{ideal membership},
}
\]

and

\[
\boxed{
\text{solving polynomial systems}.
}
\]

The choice of monomial ordering strongly affects the resulting basis and
computational cost.

%%%%%%%%%%%%%%%%%%%%%%%%%%%%%%%%%%%%%%%%%%%%%%%%%%%%%%%%%%%%
\subsection{Symbolic Linear Algebra}
\label{subsec:symbolic-linear-algebra}
%%%%%%%%%%%%%%%%%%%%%%%%%%%%%%%%%%%%%%%%%%%%%%%%%%%%%%%%%%%%

Matrices may contain symbolic entries.

For example,

\[
\boxed{
A
=
\begin{pmatrix}
a&b\\
c&d
\end{pmatrix}.
}
\]

Then its determinant is

\[
\boxed{
\det(A)
=
ad-bc.
}
\]

Symbolic linear algebra preserves exact dependence on parameters.

%%%%%%%%%%%%%%%%%%%%%%%%%%%%%%%%%%%%%%%%%%%%%%%%%%%%%%%%%%%%
\subsection{Symbolic Matrix Inverse}
\label{subsec:symbolic-matrix-inverse}
%%%%%%%%%%%%%%%%%%%%%%%%%%%%%%%%%%%%%%%%%%%%%%%%%%%%%%%%%%%%

For

\[
A
=
\begin{pmatrix}
a&b\\
c&d
\end{pmatrix},
\]

if

\[
ad-bc\neq0,
\]

then

\[
\boxed{
A^{-1}
=
\frac1{ad-bc}
\begin{pmatrix}
d&-b\\
-c&a
\end{pmatrix}.
}
\]

The nonzero determinant condition is part of the symbolic result.

%%%%%%%%%%%%%%%%%%%%%%%%%%%%%%%%%%%%%%%%%%%%%%%%%%%%%%%%%%%%
\subsection{Symbolic Gaussian Elimination}
\label{subsec:symbolic-gaussian-elimination}
%%%%%%%%%%%%%%%%%%%%%%%%%%%%%%%%%%%%%%%%%%%%%%%%%%%%%%%%%%%%

Gaussian elimination can be performed exactly with symbolic entries.

However, expressions may grow dramatically.

For example, eliminating variables from a symbolic matrix can generate
large rational expressions.

Thus exactness does not imply computational simplicity.

%%%%%%%%%%%%%%%%%%%%%%%%%%%%%%%%%%%%%%%%%%%%%%%%%%%%%%%%%%%%
\subsection{Symbolic Determinants}
\label{subsec:symbolic-determinants}
%%%%%%%%%%%%%%%%%%%%%%%%%%%%%%%%%%%%%%%%%%%%%%%%%%%%%%%%%%%%

A determinant may be expanded using cofactors:

\[
\boxed{
\det(A)
=
\sum_{j=1}^{n}
(-1)^{i+j}
a_{ij}
M_{ij}.
}
\]

But direct cofactor expansion has poor computational complexity.

Practical symbolic systems use elimination and structured algorithms when
possible.

%%%%%%%%%%%%%%%%%%%%%%%%%%%%%%%%%%%%%%%%%%%%%%%%%%%%%%%%%%%%
\subsection{Symbolic Eigenvalues}
\label{subsec:symbolic-eigenvalues}
%%%%%%%%%%%%%%%%%%%%%%%%%%%%%%%%%%%%%%%%%%%%%%%%%%%%%%%%%%%%

Eigenvalues satisfy

\[
\boxed{
\det(A-\lambda I)=0.
}
\]

For small symbolic matrices, this can produce exact parameter-dependent
eigenvalues.

For example,

\[
A
=
\begin{pmatrix}
a&0\\
0&b
\end{pmatrix}
\]

has

\[
\boxed{
\lambda_1=a,
\qquad
\lambda_2=b.
}
\]

For larger general matrices, exact symbolic eigenvalues may become
impractical.

%%%%%%%%%%%%%%%%%%%%%%%%%%%%%%%%%%%%%%%%%%%%%%%%%%%%%%%%%%%%
\subsection{Expression Swell}
\label{subsec:expression-swell}
%%%%%%%%%%%%%%%%%%%%%%%%%%%%%%%%%%%%%%%%%%%%%%%%%%%%%%%%%%%%

Expression swell occurs when intermediate symbolic expressions become much
larger than the final result.

For example, repeated expansion may produce thousands of terms that later
cancel.

Expression swell can dominate:

\[
\boxed{
\text{memory},
}
\]

and

\[
\boxed{
\text{runtime}.
}
\]

Controlling intermediate representation is therefore critical.

%%%%%%%%%%%%%%%%%%%%%%%%%%%%%%%%%%%%%%%%%%%%%%%%%%%%%%%%%%%%
\subsection{Avoiding Unnecessary Expansion}
\label{subsec:avoid-unnecessary-expansion}
%%%%%%%%%%%%%%%%%%%%%%%%%%%%%%%%%%%%%%%%%%%%%%%%%%%%%%%%%%%%

Consider

\[
(x_1+\cdots+x_n)^{20}.
\]

Keeping the expression in compact power form requires little storage.

Fully expanding it may create an enormous number of terms.

Symbolic algorithms should therefore expand only when expansion serves a
specific purpose.

%%%%%%%%%%%%%%%%%%%%%%%%%%%%%%%%%%%%%%%%%%%%%%%%%%%%%%%%%%%%
\subsection{Common Subexpression Elimination}
\label{subsec:common-subexpression-elimination}
%%%%%%%%%%%%%%%%%%%%%%%%%%%%%%%%%%%%%%%%%%%%%%%%%%%%%%%%%%%%

Suppose an expression contains

\[
x^2+y^2
\]

many times.

Define

\[
\boxed{
t=x^2+y^2.
}
\]

Then replace repeated occurrences by \(t\).

This reduces:

\[
\boxed{
\text{expression size},
}
\]

and

\[
\boxed{
\text{repeated computation}.
}
\]

Common subexpression elimination is important when symbolic formulas are
converted into numerical code.

%%%%%%%%%%%%%%%%%%%%%%%%%%%%%%%%%%%%%%%%%%%%%%%%%%%%%%%%%%%%
\subsection{Pattern Matching}
\label{subsec:symbolic-pattern-matching}
%%%%%%%%%%%%%%%%%%%%%%%%%%%%%%%%%%%%%%%%%%%%%%%%%%%%%%%%%%%%

Pattern matching identifies expressions with a specified structural form.

For example, a rule might recognize

\[
\boxed{
u+0
}
\]

and replace it by

\[
\boxed{
u.
}
\]

Another rule might recognize

\[
\boxed{
u^m u^n
}
\]

and combine powers under appropriate assumptions.

%%%%%%%%%%%%%%%%%%%%%%%%%%%%%%%%%%%%%%%%%%%%%%%%%%%%%%%%%%%%
\subsection{Rewrite Rules}
\label{subsec:symbolic-rewrite-rules}
%%%%%%%%%%%%%%%%%%%%%%%%%%%%%%%%%%%%%%%%%%%%%%%%%%%%%%%%%%%%

A rewrite rule has form

\[
\boxed{
\text{pattern}
\longrightarrow
\text{replacement}.
}
\]

Examples include:

\[
x+0
\longrightarrow
x,
\]

\[
x\cdot1
\longrightarrow
x,
\]

and

\[
\boxed{
\sin^2x+\cos^2x
\longrightarrow
1.
}
\]

Rules may be algebraic, trigonometric, logical, or domain dependent.

%%%%%%%%%%%%%%%%%%%%%%%%%%%%%%%%%%%%%%%%%%%%%%%%%%%%%%%%%%%%
\subsection{Term-Rewriting Systems}
\label{subsec:term-rewriting-systems}
%%%%%%%%%%%%%%%%%%%%%%%%%%%%%%%%%%%%%%%%%%%%%%%%%%%%%%%%%%%%

A term-rewriting system repeatedly applies transformation rules until no
more applicable rules remain or a desired form is reached.

Important questions include:

\[
\boxed{
\text{Does rewriting terminate?}
}
\]

and

\[
\boxed{
\text{Is the final result independent of rule order?}
}
\]

These lead to the concepts of termination and confluence.

%%%%%%%%%%%%%%%%%%%%%%%%%%%%%%%%%%%%%%%%%%%%%%%%%%%%%%%%%%%%
\subsection{Termination}
\label{subsec:rewriting-termination}
%%%%%%%%%%%%%%%%%%%%%%%%%%%%%%%%%%%%%%%%%%%%%%%%%%%%%%%%%%%%

A rewriting process terminates if it cannot continue indefinitely.

Poorly designed rules can create loops.

For example,

\[
a
\longrightarrow
b
\]

combined with

\[
b
\longrightarrow
a
\]

produces infinite rewriting.

Symbolic simplifiers therefore require carefully oriented rules.

%%%%%%%%%%%%%%%%%%%%%%%%%%%%%%%%%%%%%%%%%%%%%%%%%%%%%%%%%%%%
\subsection{Confluence}
\label{subsec:rewriting-confluence}
%%%%%%%%%%%%%%%%%%%%%%%%%%%%%%%%%%%%%%%%%%%%%%%%%%%%%%%%%%%%

A rewriting system is confluent if different valid sequences of
rewrites eventually lead to the same normal form.

Confluence is valuable because the final result does not depend on the
order in which rules are applied.

Not all symbolic transformation systems are confluent.

%%%%%%%%%%%%%%%%%%%%%%%%%%%%%%%%%%%%%%%%%%%%%%%%%%%%%%%%%%%%
\subsection{Assumptions}
\label{subsec:symbolic-assumptions}
%%%%%%%%%%%%%%%%%%%%%%%%%%%%%%%%%%%%%%%%%%%%%%%%%%%%%%%%%%%%

Many symbolic identities depend on assumptions.

For example,

\[
\sqrt{x^2}
\]

is not always equal to \(x\) over the real numbers.

Instead,

\[
\boxed{
\sqrt{x^2}
=
|x|.
}
\]

If the assumption

\[
x>0
\]

is known, then

\[
\boxed{
\sqrt{x^2}=x.
}
\]

%%%%%%%%%%%%%%%%%%%%%%%%%%%%%%%%%%%%%%%%%%%%%%%%%%%%%%%%%%%%
\subsection{Real and Complex Domains}
\label{subsec:real-complex-symbolic-domains}
%%%%%%%%%%%%%%%%%%%%%%%%%%%%%%%%%%%%%%%%%%%%%%%%%%%%%%%%%%%%

An identity valid over the positive real numbers may fail over all real or
complex values.

For example,

\[
\boxed{
\sqrt{ab}
=
\sqrt{a}\sqrt{b}
}
\]

is not universally valid for arbitrary complex \(a\) and \(b\) under
principal-branch conventions.

Domains and branches therefore matter in symbolic simplification.

%%%%%%%%%%%%%%%%%%%%%%%%%%%%%%%%%%%%%%%%%%%%%%%%%%%%%%%%%%%%
\subsection{Branch Cuts}
\label{subsec:symbolic-branch-cuts}
%%%%%%%%%%%%%%%%%%%%%%%%%%%%%%%%%%%%%%%%%%%%%%%%%%%%%%%%%%%%

Functions such as

\[
\log z,
\]

\[
\sqrt{z},
\]

and inverse trigonometric functions are multivalued naturally in complex
analysis.

A symbolic system usually selects principal branches.

Branch cuts specify where these branches become discontinuous.

Ignoring them can produce false symbolic identities.

%%%%%%%%%%%%%%%%%%%%%%%%%%%%%%%%%%%%%%%%%%%%%%%%%%%%%%%%%%%%
\subsection{Piecewise Expressions}
\label{subsec:symbolic-piecewise-expressions}
%%%%%%%%%%%%%%%%%%%%%%%%%%%%%%%%%%%%%%%%%%%%%%%%%%%%%%%%%%%%

Some exact mathematical results require piecewise descriptions.

For example,

\[
|x|
=
\boxed{
\begin{cases}
x,
&
x\geq0,
\\
-x,
&
x<0.
\end{cases}
}
\]

Similarly, solving inequalities or integrating expressions with absolute
values may naturally produce piecewise answers.

%%%%%%%%%%%%%%%%%%%%%%%%%%%%%%%%%%%%%%%%%%%%%%%%%%%%%%%%%%%%
\subsection{Symbolic Inequalities}
\label{subsec:symbolic-inequalities}
%%%%%%%%%%%%%%%%%%%%%%%%%%%%%%%%%%%%%%%%%%%%%%%%%%%%%%%%%%%%

Solving

\[
f(x)>0
\]

requires determining where the function has the correct sign.

For a polynomial, symbolic methods can use:

\[
\boxed{
\text{factorization},
}
\]

\[
\boxed{
\text{root isolation},
}
\]

and

\[
\boxed{
\text{sign analysis}.
}
\]

The result often consists of intervals rather than isolated points.

%%%%%%%%%%%%%%%%%%%%%%%%%%%%%%%%%%%%%%%%%%%%%%%%%%%%%%%%%%%%
\subsection{Worked Example: Symbolic Inequality}
\label{subsec:worked-symbolic-inequality}
%%%%%%%%%%%%%%%%%%%%%%%%%%%%%%%%%%%%%%%%%%%%%%%%%%%%%%%%%%%%

\begin{workedexamplebox}{Solving a Polynomial Inequality}

Solve

\[
x^2-5x+6>0.
\]

Factor:

\[
(x-2)(x-3)>0.
\]

The product is positive when both factors have the same sign.

Thus,

\[
x<2
\]

or

\[
x>3.
\]

\begin{answerbox}
\[
\boxed{
(-\infty,2)
\cup
(3,\infty).
}
\]
\end{answerbox}

\end{workedexamplebox}

%%%%%%%%%%%%%%%%%%%%%%%%%%%%%%%%%%%%%%%%%%%%%%%%%%%%%%%%%%%%
\subsection{Symbolic Summation}
\label{subsec:symbolic-summation}
%%%%%%%%%%%%%%%%%%%%%%%%%%%%%%%%%%%%%%%%%%%%%%%%%%%%%%%%%%%%

Symbolic summation seeks closed forms for expressions such as

\[
\boxed{
\sum_{k=1}^{n}
f(k).
}
\]

For example,

\[
\boxed{
\sum_{k=1}^{n}
k
=
\frac{
n(n+1)
}{
2
}.
}
\]

Likewise,

\[
\boxed{
\sum_{k=1}^{n}
k^2
=
\frac{
n(n+1)(2n+1)
}{
6
}.
}
\]

%%%%%%%%%%%%%%%%%%%%%%%%%%%%%%%%%%%%%%%%%%%%%%%%%%%%%%%%%%%%
\subsection{Geometric Series}
\label{subsec:symbolic-geometric-series}
%%%%%%%%%%%%%%%%%%%%%%%%%%%%%%%%%%%%%%%%%%%%%%%%%%%%%%%%%%%%

For

\[
r\neq1,
\]

\[
\boxed{
\sum_{k=0}^{n}
r^k
=
\frac{
1-r^{n+1}
}{
1-r
}.
}
\]

For

\[
|r|<1,
\]

letting \(n\to\infty\) gives

\[
\boxed{
\sum_{k=0}^{\infty}
r^k
=
\frac1{1-r}.
}
\]

Symbolic summation can manipulate both finite and infinite sums.

%%%%%%%%%%%%%%%%%%%%%%%%%%%%%%%%%%%%%%%%%%%%%%%%%%%%%%%%%%%%
\subsection{Symbolic Products}
\label{subsec:symbolic-products}
%%%%%%%%%%%%%%%%%%%%%%%%%%%%%%%%%%%%%%%%%%%%%%%%%%%%%%%%%%%%

Finite products may also admit closed forms.

For example,

\[
\boxed{
\prod_{k=1}^{n}
k
=
n!.
}
\]

Another example is

\[
\boxed{
\prod_{k=1}^{n}
\frac{k+1}{k}
=
n+1.
}
\]

The second product telescopes.

%%%%%%%%%%%%%%%%%%%%%%%%%%%%%%%%%%%%%%%%%%%%%%%%%%%%%%%%%%%%
\subsection{Telescoping Structures}
\label{subsec:symbolic-telescoping}
%%%%%%%%%%%%%%%%%%%%%%%%%%%%%%%%%%%%%%%%%%%%%%%%%%%%%%%%%%%%

A sum telescopes when consecutive terms cancel.

For example,

\[
\sum_{k=1}^{n}
\left(
\frac1k
-
\frac1{k+1}
\right)
\]

becomes

\[
\boxed{
1-\frac1{n+1}.
}
\]

Recognizing telescoping structure is an important symbolic transformation.

%%%%%%%%%%%%%%%%%%%%%%%%%%%%%%%%%%%%%%%%%%%%%%%%%%%%%%%%%%%%
\subsection{Recurrence Relations}
\label{subsec:symbolic-recurrence-relations}
%%%%%%%%%%%%%%%%%%%%%%%%%%%%%%%%%%%%%%%%%%%%%%%%%%%%%%%%%%%%

A recurrence defines sequence values using previous values.

For example,

\[
\boxed{
a_{n+1}
=
ra_n.
}
\]

With

\[
a_0=C,
\]

the symbolic solution is

\[
\boxed{
a_n
=
Cr^n.
}
\]

More complicated linear recurrences can be solved using characteristic
polynomials.

%%%%%%%%%%%%%%%%%%%%%%%%%%%%%%%%%%%%%%%%%%%%%%%%%%%%%%%%%%%%
\subsection{Linear Homogeneous Recurrences}
\label{subsec:linear-homogeneous-recurrences-symbolic}
%%%%%%%%%%%%%%%%%%%%%%%%%%%%%%%%%%%%%%%%%%%%%%%%%%%%%%%%%%%%

Consider

\[
\boxed{
a_{n+2}
-
3a_{n+1}
+
2a_n
=
0.
}
\]

Try

\[
a_n=r^n.
\]

Then

\[
r^2-3r+2=0.
\]

Factoring,

\[
(r-1)(r-2)=0.
\]

Therefore,

\[
\boxed{
a_n
=
C_1
+
C_2 2^n.
}
\]

Initial values determine \(C_1\) and \(C_2\).

%%%%%%%%%%%%%%%%%%%%%%%%%%%%%%%%%%%%%%%%%%%%%%%%%%%%%%%%%%%%
\subsection{Symbolic Differential Equations}
\label{subsec:symbolic-differential-equations}
%%%%%%%%%%%%%%%%%%%%%%%%%%%%%%%%%%%%%%%%%%%%%%%%%%%%%%%%%%%%

Some differential equations admit exact symbolic solutions.

For example,

\[
y'=ky
\]

has solution

\[
\boxed{
y(t)
=
Ce^{kt}.
}
\]

With initial condition

\[
y(0)=y_0,
\]

we obtain

\[
\boxed{
y(t)
=
y_0e^{kt}.
}
\]

%%%%%%%%%%%%%%%%%%%%%%%%%%%%%%%%%%%%%%%%%%%%%%%%%%%%%%%%%%%%
\subsection{Symbolic Separation of Variables}
\label{subsec:symbolic-separation-variables}
%%%%%%%%%%%%%%%%%%%%%%%%%%%%%%%%%%%%%%%%%%%%%%%%%%%%%%%%%%%%

For

\[
y'=g(t)h(y),
\]

formally separate:

\[
\boxed{
\frac{dy}{h(y)}
=
g(t)\,dt.
}
\]

Integrating gives

\[
\boxed{
\int
\frac{dy}{h(y)}
=
\int
g(t)\,dt.
}
\]

The resulting implicit relation may then be solved for \(y\) if possible.

%%%%%%%%%%%%%%%%%%%%%%%%%%%%%%%%%%%%%%%%%%%%%%%%%%%%%%%%%%%%
\subsection{Symbolic Linear ODEs}
\label{subsec:symbolic-linear-odes}
%%%%%%%%%%%%%%%%%%%%%%%%%%%%%%%%%%%%%%%%%%%%%%%%%%%%%%%%%%%%

A first-order linear equation has form

\[
\boxed{
y'
+
p(x)y
=
q(x).
}
\]

Its integrating factor is

\[
\boxed{
\mu(x)
=
e^{\int p(x)\,dx}.
}
\]

Then

\[
\boxed{
\frac{d}{dx}
[
\mu(x)y(x)
]
=
\mu(x)q(x).
}
\]

Symbolic systems can automate this procedure for many input functions.

%%%%%%%%%%%%%%%%%%%%%%%%%%%%%%%%%%%%%%%%%%%%%%%%%%%%%%%%%%%%
\subsection{Symbolic Laplace Transforms}
\label{subsec:symbolic-laplace-transforms}
%%%%%%%%%%%%%%%%%%%%%%%%%%%%%%%%%%%%%%%%%%%%%%%%%%%%%%%%%%%%

The Laplace transform is

\[
\boxed{
\mathcal{L}\{f(t)\}(s)
=
\int_0^\infty
e^{-st}f(t)\,dt.
}
\]

For example,

\[
\boxed{
\mathcal{L}\{1\}
=
\frac1s,
}
\]

and

\[
\boxed{
\mathcal{L}\{e^{at}\}
=
\frac1{s-a}.
}
\]

Symbolic transforms convert many differential equations into algebraic
equations.

%%%%%%%%%%%%%%%%%%%%%%%%%%%%%%%%%%%%%%%%%%%%%%%%%%%%%%%%%%%%
\subsection{Symbolic Fourier Transforms}
\label{subsec:symbolic-fourier-transforms}
%%%%%%%%%%%%%%%%%%%%%%%%%%%%%%%%%%%%%%%%%%%%%%%%%%%%%%%%%%%%

A Fourier transform can be written

\[
\boxed{
\widehat{f}(\omega)
=
\int_{-\infty}^{\infty}
f(x)e^{-\mathrm{i}\omega x}\,dx,
}
\]

under one common convention.

Symbolic transform systems require careful handling of normalization,
domains, distributions, and convergence assumptions.

%%%%%%%%%%%%%%%%%%%%%%%%%%%%%%%%%%%%%%%%%%%%%%%%%%%%%%%%%%%%
\subsection{Symbolic Optimization}
\label{subsec:symbolic-optimization}
%%%%%%%%%%%%%%%%%%%%%%%%%%%%%%%%%%%%%%%%%%%%%%%%%%%%%%%%%%%%

For simple objectives, symbolic differentiation can identify stationary
points exactly.

Given

\[
f(x),
\]

solve

\[
\boxed{
f'(x)=0.
}
\]

Then classify candidates using

\[
\boxed{
f''(x)
}
\]

or direct comparison.

For multivariable problems, solve

\[
\boxed{
\nabla f(x)=0.
}
\]

%%%%%%%%%%%%%%%%%%%%%%%%%%%%%%%%%%%%%%%%%%%%%%%%%%%%%%%%%%%%
\subsection{Worked Example: Symbolic Optimization}
\label{subsec:worked-symbolic-optimization}
%%%%%%%%%%%%%%%%%%%%%%%%%%%%%%%%%%%%%%%%%%%%%%%%%%%%%%%%%%%%

\begin{workedexamplebox}{Exact Minimizer of a Quadratic}

Consider

\[
f(x)
=
3x^2-12x+7.
\]

Differentiate:

\[
f'(x)
=
6x-12.
\]

Set

\[
6x-12=0.
\]

Thus,

\[
x=2.
\]

Since

\[
f''(x)=6>0,
\]

this is the unique minimizer.

\begin{answerbox}
\[
\boxed{
x^*=2.
}
\]
\end{answerbox}

\end{workedexamplebox}

%%%%%%%%%%%%%%%%%%%%%%%%%%%%%%%%%%%%%%%%%%%%%%%%%%%%%%%%%%%%
\subsection{Symbolic Jacobians}
\label{subsec:symbolic-jacobians}
%%%%%%%%%%%%%%%%%%%%%%%%%%%%%%%%%%%%%%%%%%%%%%%%%%%%%%%%%%%%

For vector function

\[
F
=
\begin{pmatrix}
f_1\\
\vdots\\
f_m
\end{pmatrix},
\]

the Jacobian is

\[
\boxed{
J_F
=
\left[
\frac{
\partial f_i
}{
\partial x_j
}
\right].
}
\]

Symbolic Jacobians are useful in:

\[
\boxed{
\text{Newton methods},
}
\]

\[
\boxed{
\text{dynamical systems},
}
\]

and

\[
\boxed{
\text{implicit differential-equation solvers}.
}
\]

%%%%%%%%%%%%%%%%%%%%%%%%%%%%%%%%%%%%%%%%%%%%%%%%%%%%%%%%%%%%
\subsection{Symbolic Hessians}
\label{subsec:symbolic-hessians}
%%%%%%%%%%%%%%%%%%%%%%%%%%%%%%%%%%%%%%%%%%%%%%%%%%%%%%%%%%%%

For scalar function \(f\), the Hessian is

\[
\boxed{
H_f
=
\left[
\frac{
\partial^2f
}{
\partial x_i\partial x_j
}
\right].
}
\]

Symbolically computing the Hessian provides exact curvature expressions for
optimization and local analysis.

%%%%%%%%%%%%%%%%%%%%%%%%%%%%%%%%%%%%%%%%%%%%%%%%%%%%%%%%%%%%
\subsection{Code Generation}
\label{subsec:symbolic-code-generation}
%%%%%%%%%%%%%%%%%%%%%%%%%%%%%%%%%%%%%%%%%%%%%%%%%%%%%%%%%%%%

A symbolic expression can be transformed into executable numerical code.

For example,

\[
f(x,y)
=
x^2
+
2xy
+
y^2
\]

may be converted into a numerical evaluation routine.

Before code generation, a symbolic system may perform:

\[
\boxed{
\text{simplification},
}
\]

\[
\boxed{
\text{common-subexpression elimination},
}
\]

and

\[
\boxed{
\text{algebraic optimization}.
}
\]

%%%%%%%%%%%%%%%%%%%%%%%%%%%%%%%%%%%%%%%%%%%%%%%%%%%%%%%%%%%%
\subsection{Symbolic-Numeric Computation}
\label{subsec:symbolic-numeric-computation}
%%%%%%%%%%%%%%%%%%%%%%%%%%%%%%%%%%%%%%%%%%%%%%%%%%%%%%%%%%%%

Many practical algorithms combine symbolic and numerical methods.

A typical workflow is:

\[
\boxed{
\text{derive symbolically}
\rightarrow
\text{simplify}
\rightarrow
\text{evaluate numerically}.
}
\]

For example, one may derive an exact Jacobian symbolically and then
evaluate it numerically inside Newton's method.

%%%%%%%%%%%%%%%%%%%%%%%%%%%%%%%%%%%%%%%%%%%%%%%%%%%%%%%%%%%%
\subsection{Exact Preprocessing}
\label{subsec:exact-preprocessing}
%%%%%%%%%%%%%%%%%%%%%%%%%%%%%%%%%%%%%%%%%%%%%%%%%%%%%%%%%%%%

Symbolic preprocessing can simplify a problem before numerical solution.

Examples include:

\[
\boxed{
\text{eliminating variables},
}
\]

\[
\boxed{
\text{factoring expressions},
}
\]

\[
\boxed{
\text{deriving exact gradients},
}
\]

and

\[
\boxed{
\text{nondimensionalization}.
}
\]

This can improve numerical efficiency and reliability.

%%%%%%%%%%%%%%%%%%%%%%%%%%%%%%%%%%%%%%%%%%%%%%%%%%%%%%%%%%%%
\subsection{Numeric Evaluation of Symbolic Expressions}
\label{subsec:numeric-evaluation-symbolic}
%%%%%%%%%%%%%%%%%%%%%%%%%%%%%%%%%%%%%%%%%%%%%%%%%%%%%%%%%%%%

After obtaining symbolic expression

\[
F(x),
\]

one may substitute

\[
x=x_0
\]

and evaluate numerically.

Precision can be chosen independently of the original symbolic expression.

For example,

\[
\sqrt{2}
\]

may be evaluated to:

\[
\boxed{
10,
\quad
50,
\quad
100,
}
\]

or more decimal digits.

%%%%%%%%%%%%%%%%%%%%%%%%%%%%%%%%%%%%%%%%%%%%%%%%%%%%%%%%%%%%
\subsection{Symbolic Verification}
\label{subsec:symbolic-verification}
%%%%%%%%%%%%%%%%%%%%%%%%%%%%%%%%%%%%%%%%%%%%%%%%%%%%%%%%%%%%

Symbolic computation can verify identities by simplifying their
difference.

To test

\[
f(x)=g(x),
\]

form

\[
\boxed{
h(x)
=
f(x)-g(x).
}
\]

If symbolic simplification proves

\[
\boxed{
h(x)=0
}
\]

under stated assumptions, the identity is verified.

%%%%%%%%%%%%%%%%%%%%%%%%%%%%%%%%%%%%%%%%%%%%%%%%%%%%%%%%%%%%
\subsection{Worked Example: Identity Verification}
\label{subsec:worked-symbolic-identity}
%%%%%%%%%%%%%%%%%%%%%%%%%%%%%%%%%%%%%%%%%%%%%%%%%%%%%%%%%%%%

\begin{workedexamplebox}{Verifying a Trigonometric Identity}

Consider

\[
1-\cos^2x
\]

and

\[
\sin^2x.
\]

Subtract:

\[
1-\cos^2x-\sin^2x.
\]

Using

\[
\sin^2x+\cos^2x=1,
\]

the expression simplifies to

\[
0.
\]

\begin{answerbox}
\[
\boxed{
1-\cos^2x
=
\sin^2x.
}
\]
\end{answerbox}

\end{workedexamplebox}

%%%%%%%%%%%%%%%%%%%%%%%%%%%%%%%%%%%%%%%%%%%%%%%%%%%%%%%%%%%%
\subsection{Symbolic Proof Assistance}
\label{subsec:symbolic-proof-assistance}
%%%%%%%%%%%%%%%%%%%%%%%%%%%%%%%%%%%%%%%%%%%%%%%%%%%%%%%%%%%%

Symbolic systems can assist mathematical proof by performing exact
subcalculations.

Examples include:

\[
\boxed{
\text{polynomial identity checking},
}
\]

\[
\boxed{
\text{factorization},
}
\]

\[
\boxed{
\text{elimination},
}
\]

and

\[
\boxed{
\text{exact differentiation}.
}
\]

However, a computational output is not automatically a human-readable
proof.

The mathematical assumptions and algorithmic guarantees still matter.

%%%%%%%%%%%%%%%%%%%%%%%%%%%%%%%%%%%%%%%%%%%%%%%%%%%%%%%%%%%%
\subsection{Symbolic Computation and Theorem Proving}
\label{subsec:symbolic-computation-theorem-proving}
%%%%%%%%%%%%%%%%%%%%%%%%%%%%%%%%%%%%%%%%%%%%%%%%%%%%%%%%%%%%

Symbolic algebra manipulates mathematical expressions.

Automated theorem proving manipulates logical statements and proof rules.

The two fields overlap when algebraic identities or inequalities are used
inside formal proofs.

However,

\[
\boxed{
\text{computer algebra}
\neq
\text{formal proof system}.
}
\]

A symbolic result may rely on algorithms rather than a proof object checked
by a logical kernel.

%%%%%%%%%%%%%%%%%%%%%%%%%%%%%%%%%%%%%%%%%%%%%%%%%%%%%%%%%%%%
\subsection{Computational Complexity}
\label{subsec:symbolic-complexity}
%%%%%%%%%%%%%%%%%%%%%%%%%%%%%%%%%%%%%%%%%%%%%%%%%%%%%%%%%%%%

Symbolic algorithms can have dramatically different complexity from
numerical algorithms.

A symbolic expression may grow exponentially in size even when the final
answer is short.

Costs therefore depend on:

\[
\boxed{
\text{number of terms},
}
\]

\[
\boxed{
\text{integer digit length},
}
\]

\[
\boxed{
\text{polynomial degree},
}
\]

and

\[
\boxed{
\text{number of variables}.
}
\]

%%%%%%%%%%%%%%%%%%%%%%%%%%%%%%%%%%%%%%%%%%%%%%%%%%%%%%%%%%%%
\subsection{Bit Complexity}
\label{subsec:symbolic-bit-complexity}
%%%%%%%%%%%%%%%%%%%%%%%%%%%%%%%%%%%%%%%%%%%%%%%%%%%%%%%%%%%%

When integers become large, counting arithmetic operations alone is not
enough.

For example, multiplying two integers with millions of digits costs more
than multiplying two machine-size integers.

Bit complexity measures work in terms of the number of bits manipulated.

This is especially important in exact arithmetic.

%%%%%%%%%%%%%%%%%%%%%%%%%%%%%%%%%%%%%%%%%%%%%%%%%%%%%%%%%%%%
\subsection{Coefficient Growth}
\label{subsec:coefficient-growth-symbolic}
%%%%%%%%%%%%%%%%%%%%%%%%%%%%%%%%%%%%%%%%%%%%%%%%%%%%%%%%%%%%

Exact polynomial operations can produce coefficients much larger than
those appearing in the input.

This coefficient growth can cause substantial memory and runtime cost.

Modular algorithms often reduce this problem by computing over smaller
finite fields and reconstructing exact results afterward.

%%%%%%%%%%%%%%%%%%%%%%%%%%%%%%%%%%%%%%%%%%%%%%%%%%%%%%%%%%%%
\subsection{Modular Computation Preview}
\label{subsec:modular-computation-preview}
%%%%%%%%%%%%%%%%%%%%%%%%%%%%%%%%%%%%%%%%%%%%%%%%%%%%%%%%%%%%

Instead of performing a difficult integer computation directly, reduce the
problem modulo several primes:

\[
\boxed{
p_1,
p_2,
\ldots,p_k.
}
\]

Perform smaller calculations modulo each prime.

Then reconstruct the integer result using tools such as the Chinese
Remainder Theorem.

This is a powerful strategy in computer algebra.

%%%%%%%%%%%%%%%%%%%%%%%%%%%%%%%%%%%%%%%%%%%%%%%%%%%%%%%%%%%%
\subsection{Evaluation and Interpolation}
\label{subsec:symbolic-evaluation-interpolation}
%%%%%%%%%%%%%%%%%%%%%%%%%%%%%%%%%%%%%%%%%%%%%%%%%%%%%%%%%%%%

Another symbolic strategy evaluates a polynomial or expression at selected
points, performs easier computations on the values, and reconstructs the
result by interpolation.

Conceptually,

\[
\boxed{
\text{symbolic object}
\rightarrow
\text{evaluations}
\rightarrow
\text{computation}
\rightarrow
\text{interpolation}.
}
\]

This can avoid intermediate expression swell.

%%%%%%%%%%%%%%%%%%%%%%%%%%%%%%%%%%%%%%%%%%%%%%%%%%%%%%%%%%%%
\subsection{Expression Hashing and Sharing}
\label{subsec:expression-sharing}
%%%%%%%%%%%%%%%%%%%%%%%%%%%%%%%%%%%%%%%%%%%%%%%%%%%%%%%%%%%%

If identical subexpressions appear repeatedly, a symbolic system may store
them once and let several expressions reference the same object.

This creates a directed acyclic graph rather than a pure tree.

Sharing reduces:

\[
\boxed{
\text{memory duplication}
}
\]

and sometimes

\[
\boxed{
\text{repeated computation}.
}
\]

%%%%%%%%%%%%%%%%%%%%%%%%%%%%%%%%%%%%%%%%%%%%%%%%%%%%%%%%%%%%
\subsection{Lazy Evaluation}
\label{subsec:symbolic-lazy-evaluation}
%%%%%%%%%%%%%%%%%%%%%%%%%%%%%%%%%%%%%%%%%%%%%%%%%%%%%%%%%%%%

Some transformations are expensive and may never be needed.

Lazy evaluation delays computation until the result is required.

For example, a huge product may remain factored until coefficient access
or expansion is requested.

This can prevent unnecessary symbolic work.

%%%%%%%%%%%%%%%%%%%%%%%%%%%%%%%%%%%%%%%%%%%%%%%%%%%%%%%%%%%%
\subsection{Simplification Heuristics}
\label{subsec:simplification-heuristics}
%%%%%%%%%%%%%%%%%%%%%%%%%%%%%%%%%%%%%%%%%%%%%%%%%%%%%%%%%%%%

Because no single measure of expression simplicity works in all contexts,
symbolic systems often use heuristics.

Possible criteria include:

\[
\boxed{
\text{number of operations},
}
\]

\[
\boxed{
\text{number of terms},
}
\]

\[
\boxed{
\text{expression depth},
}
\]

and

\[
\boxed{
\text{preferred mathematical functions}.
}
\]

Different simplifiers may therefore return different but equivalent forms.

%%%%%%%%%%%%%%%%%%%%%%%%%%%%%%%%%%%%%%%%%%%%%%%%%%%%%%%%%%%%
\subsection{Symbolic Computation Workflow}
\label{subsec:symbolic-computation-workflow}
%%%%%%%%%%%%%%%%%%%%%%%%%%%%%%%%%%%%%%%%%%%%%%%%%%%%%%%%%%%%

A practical workflow is:

\begin{enumerate}
    \item define symbols and their domains;
    \item preserve exact constants whenever exactness matters;
    \item construct the expression structurally;
    \item choose the desired algebraic form;
    \item expand, factor, collect, or cancel only when useful;
    \item state assumptions explicitly;
    \item perform differentiation, integration, solving, or elimination;
    \item simplify the result;
    \item check excluded domain points;
    \item check branches of multivalued functions;
    \item substitute the result back into the original equation when
          possible;
    \item compare derivatives to verify antiderivatives;
    \item use numerical evaluation as an independent check;
    \item avoid unnecessary expansion;
    \item control expression swell;
    \item convert symbolic expressions to numerical code when needed;
    \item document assumptions and transformations.
\end{enumerate}

%%%%%%%%%%%%%%%%%%%%%%%%%%%%%%%%%%%%%%%%%%%%%%%%%%%%%%%%%%%%
\subsection{Choosing Symbolic or Numerical Methods}
\label{subsec:choosing-symbolic-or-numerical}
%%%%%%%%%%%%%%%%%%%%%%%%%%%%%%%%%%%%%%%%%%%%%%%%%%%%%%%%%%%%

Use symbolic methods when:

\[
\boxed{
\text{exact structure matters},
}
\]

\[
\boxed{
\text{parameters must remain symbolic},
}
\]

or

\[
\boxed{
\text{an exact identity is desired}.
}
\]

Use numerical methods when:

\[
\boxed{
\text{the exact form is unavailable},
}
\]

\[
\boxed{
\text{expressions become too large},
}
\]

or

\[
\boxed{
\text{only approximate numerical values are needed}.
}
\]

Often the best approach is hybrid.

%%%%%%%%%%%%%%%%%%%%%%%%%%%%%%%%%%%%%%%%%%%%%%%%%%%%%%%%%%%%
\subsection{Common Errors}
\label{subsec:symbolic-computation-common-errors}
%%%%%%%%%%%%%%%%%%%%%%%%%%%%%%%%%%%%%%%%%%%%%%%%%%%%%%%%%%%%

\subsubsection{Treating Approximate Input as Exact}

Entering

\[
0.333333
\]

instead of

\[
\frac13
\]

can change an exact symbolic calculation into an approximate one.

Use exact quantities when exactness is intended.

%%%%%%%%%%%%%%%%%%%%%%%%%%%%%%%%%%%%%%%%%%%%%%%%%%%%%%%%%%%%
\subsubsection{Assuming Algebraic Cancellation Preserves the Original Domain}

The simplification

\[
\frac{x^2-1}{x-1}
=
x+1
\]

is valid only for

\[
x\neq1.
\]

The simplified expression alone does not preserve the original excluded
point.

%%%%%%%%%%%%%%%%%%%%%%%%%%%%%%%%%%%%%%%%%%%%%%%%%%%%%%%%%%%%
\subsubsection{Assuming \(\sqrt{x^2}=x\)}

Over the real numbers,

\[
\boxed{
\sqrt{x^2}=|x|.
}
\]

Additional assumptions are required to replace this by \(x\).

%%%%%%%%%%%%%%%%%%%%%%%%%%%%%%%%%%%%%%%%%%%%%%%%%%%%%%%%%%%%
\subsubsection{Ignoring Complex Branches}

Identities involving logarithms, powers, square roots, and inverse
functions may change across branches in the complex plane.

Do not apply real-variable simplifications blindly.

%%%%%%%%%%%%%%%%%%%%%%%%%%%%%%%%%%%%%%%%%%%%%%%%%%%%%%%%%%%%
\subsubsection{Assuming Simplify Has One Unique Answer}

An expanded expression, factored expression, and partial-fraction
expression may all be equally valid.

The useful form depends on the task.

%%%%%%%%%%%%%%%%%%%%%%%%%%%%%%%%%%%%%%%%%%%%%%%%%%%%%%%%%%%%
\subsubsection{Expanding Everything}

Full expansion may create enormous intermediate expressions.

Preserve factored or structured forms until expansion is needed.

%%%%%%%%%%%%%%%%%%%%%%%%%%%%%%%%%%%%%%%%%%%%%%%%%%%%%%%%%%%%
\subsubsection{Factoring Over the Wrong Domain}

A polynomial irreducible over

\[
\mathbb{Q}
\]

may factor over

\[
\mathbb{R}
\]

or

\[
\mathbb{C}.
\]

Always specify or understand the coefficient domain.

%%%%%%%%%%%%%%%%%%%%%%%%%%%%%%%%%%%%%%%%%%%%%%%%%%%%%%%%%%%%
\subsubsection{Assuming Every Integral Has an Elementary Closed Form}

Many elementary integrands do not possess elementary antiderivatives.

A special function or unevaluated integral may be the correct symbolic
answer.

%%%%%%%%%%%%%%%%%%%%%%%%%%%%%%%%%%%%%%%%%%%%%%%%%%%%%%%%%%%%
\subsubsection{Assuming Every Polynomial Has a Radical Formula}

General degree-five and higher polynomials may not be solvable by radicals.

Exact implicit root representations remain possible.

%%%%%%%%%%%%%%%%%%%%%%%%%%%%%%%%%%%%%%%%%%%%%%%%%%%%%%%%%%%%
\subsubsection{Confusing Symbolic Equality With Numerical Agreement}

Two expressions can agree numerically at many sampled points without being
identically equal.

Likewise, numerical disagreement can arise from singularities or branch
choices rather than algebraic inequality.

%%%%%%%%%%%%%%%%%%%%%%%%%%%%%%%%%%%%%%%%%%%%%%%%%%%%%%%%%%%%
\subsubsection{Ignoring Assumptions in Inequality Solving}

Multiplying an inequality by an expression of unknown sign can reverse the
inequality.

Symbolic solvers must reason about sign conditions.

%%%%%%%%%%%%%%%%%%%%%%%%%%%%%%%%%%%%%%%%%%%%%%%%%%%%%%%%%%%%
\subsubsection{Ignoring Expression Swell}

An algorithm that is mathematically straightforward may be computationally
impractical because intermediate expressions become too large.

Representation matters.

%%%%%%%%%%%%%%%%%%%%%%%%%%%%%%%%%%%%%%%%%%%%%%%%%%%%%%%%%%%%
\subsubsection{Using Symbolic Matrix Inversion for Large Systems}

Exact symbolic inversion can produce enormous expressions.

Factorization, elimination, or numerical solving is usually preferable for
large parameterized systems.

%%%%%%%%%%%%%%%%%%%%%%%%%%%%%%%%%%%%%%%%%%%%%%%%%%%%%%%%%%%%
\subsubsection{Using Symbolic Eigenvalue Formulas for Large General Matrices}

Characteristic polynomials and radical expressions rapidly become
impractical.

Numerical eigenvalue methods are generally preferable for large matrices.

%%%%%%%%%%%%%%%%%%%%%%%%%%%%%%%%%%%%%%%%%%%%%%%%%%%%%%%%%%%%
\subsubsection{Trusting a Symbolic Result Without Substitution or Differentiation Checks}

Whenever practical:

\[
\boxed{
\text{substitute solutions back},
}
\]

or

\[
\boxed{
\text{differentiate antiderivatives}.
}
\]

Independent checking helps detect assumptions or branch issues.

%%%%%%%%%%%%%%%%%%%%%%%%%%%%%%%%%%%%%%%%%%%%%%%%%%%%%%%%%%%%
\subsubsection{Confusing Symbolic Computation With Formal Proof}

A symbolic transformation may be mathematically correct without generating
a proof object in a formal logical system.

The distinction matters in theorem verification.

%%%%%%%%%%%%%%%%%%%%%%%%%%%%%%%%%%%%%%%%%%%%%%%%%%%%%%%%%%%%
\subsection{Applications}
\label{subsec:symbolic-computation-applications}
%%%%%%%%%%%%%%%%%%%%%%%%%%%%%%%%%%%%%%%%%%%%%%%%%%%%%%%%%%%%

\begin{applicationbox}

\textbf{Algebra.}
Symbolic systems expand, factor, simplify, solve polynomial equations, and
perform exact elimination.

\medskip

\textbf{Calculus.}
Differentiation, integration, limits, and series can often be manipulated
exactly.

\medskip

\textbf{Differential Equations.}
Exact solutions, transforms, recurrence relations, and operator
manipulations can be automated for important equation classes.

\medskip

\textbf{Optimization.}
Symbolic gradients, Hessians, Jacobians, and KKT equations can be derived
before numerical solution.

\medskip

\textbf{Geometry.}
Exact coordinate transformations, intersection equations, and polynomial
constraints can be manipulated symbolically.

\medskip

\textbf{Number Theory.}
Exact integer arithmetic, modular computation, polynomial factorization,
and algebraic numbers are central symbolic tasks.

\medskip

\textbf{Scientific Computing.}
Symbolic preprocessing can derive efficient formulas that are later
evaluated numerically.

\medskip

\textbf{Mathematical Research.}
Computer algebra can assist conjecture formation, exact experimentation,
identity checking, and derivation of complicated formulas.

\end{applicationbox}

%%%%%%%%%%%%%%%%%%%%%%%%%%%%%%%%%%%%%%%%%%%%%%%%%%%%%%%%%%%%
\subsection{Exercises}
\label{subsec:symbolic-computation-exercises}
%%%%%%%%%%%%%%%%%%%%%%%%%%%%%%%%%%%%%%%%%%%%%%%%%%%%%%%%%%%%

\begin{exercise}[1]
Define symbolic computation.
\end{exercise}

\begin{exercise}[1]
Distinguish symbolic and numerical computation.
\end{exercise}

\begin{exercise}[1]
Define a canonical form.
\end{exercise}

\begin{exercise}[1]
Define expression swell.
\end{exercise}

\begin{exercise}[1]
Define a rewrite rule.
\end{exercise}

\begin{exercise}[1]
Define exact rational arithmetic.
\end{exercise}

\begin{exercise}[1]
Define symbolic differentiation.
\end{exercise}

\begin{exercise}[1]
Define symbolic integration.
\end{exercise}

\begin{exercise}[1]
Define a polynomial gcd.
\end{exercise}

\begin{exercise}[1]
Define common subexpression elimination.
\end{exercise}

\begin{exercise}[2]
Expand

\[
(x+3)(x-4).
\]
\end{exercise}

\begin{exercise}[2]
Factor

\[
x^2-7x+12.
\]
\end{exercise}

\begin{exercise}[2]
Collect terms in

\[
ax^2+bx+cx^2+dx+e.
\]
\end{exercise}

\begin{exercise}[2]
Reduce

\[
\frac{18}{24}
\]

to lowest terms.
\end{exercise}

\begin{exercise}[3]
Explain why

\[
\frac{x^2-4}{x-2}
\]

and

\[
x+2
\]

do not have exactly the same domain.
\end{exercise}

\begin{exercise}[2]
Add two polynomials using coefficient representations.
\end{exercise}

\begin{exercise}[2]
Multiply two polynomials symbolically.
\end{exercise}

\begin{exercise}[3]
Perform polynomial long division.
\end{exercise}

\begin{exercise}[3]
Use the Euclidean algorithm to find a polynomial gcd.
\end{exercise}

\begin{exercise}[3]
Use

\[
\gcd(f,f')
\]

to determine whether a polynomial has repeated roots.
\end{exercise}

\begin{exercise}[3]
Factor a polynomial over

\[
\mathbb{Q},
\]

then discuss whether it factors further over

\[
\mathbb{R}.
\]
\end{exercise}

\begin{exercise}[2]
Simplify a rational function by canceling polynomial gcds.
\end{exercise}

\begin{exercise}[3]
Perform partial-fraction decomposition.
\end{exercise}

\begin{exercise}[2]
Differentiate

\[
x^3e^x.
\]
\end{exercise}

\begin{exercise}[3]
Differentiate

\[
\frac{\sin(x^2)}{1+x^2}.
\]
\end{exercise}

\begin{exercise}[3]
Construct the expression tree for

\[
x^2\sin(x+1).
\]
\end{exercise}

\begin{exercise}[3]
Describe how recursive differentiation acts on the tree.
\end{exercise}

\begin{exercise}[2]
Integrate

\[
x^4.
\]
\end{exercise}

\begin{exercise}[3]
Use substitution to integrate

\[
2xe^{x^2}.
\]
\end{exercise}

\begin{exercise}[3]
Use integration by parts to integrate

\[
x\cos x.
\]
\end{exercise}

\begin{exercise}[3]
Give an example of a nonelementary antiderivative.
\end{exercise}

\begin{exercise}[2]
Evaluate symbolically

\[
\lim_{x\to2}
\frac{x^2-4}{x-2}.
\]
\end{exercise}

\begin{exercise}[3]
Find the cubic Taylor polynomial of

\[
\sin x
\]

about \(0\).
\end{exercise}

\begin{exercise}[3]
Find a Laurent expansion of

\[
\frac{e^x}{x}
\]

near \(x=0\).
\end{exercise}

\begin{exercise}[3]
Explain the difference between a convergent series and an asymptotic
series.
\end{exercise}

\begin{exercise}[2]
Solve a quadratic equation symbolically.
\end{exercise}

\begin{exercise}[3]
Solve a two-equation polynomial system by substitution.
\end{exercise}

\begin{exercise}[3]
Explain why general quintics cannot always be solved by radicals.
\end{exercise}

\begin{exercise}[3]
Explain how an exact root can be represented without radicals.
\end{exercise}

\begin{exercise}[3]
Describe the purpose of a resultant.
\end{exercise}

\begin{exercise}[3]
Explain conceptually how a Gröbner basis aids elimination.
\end{exercise}

\begin{exercise}[2]
Compute symbolically the determinant of

\[
\begin{pmatrix}
a&b\\
c&d
\end{pmatrix}.
\]
\end{exercise}

\begin{exercise}[3]
Compute the symbolic inverse of a \(2\times2\) matrix.
\end{exercise}

\begin{exercise}[3]
State the condition under which that inverse exists.
\end{exercise}

\begin{exercise}[3]
Compute the symbolic eigenvalues of a triangular matrix.
\end{exercise}

\begin{exercise}[3]
Give an example of expression swell during symbolic elimination.
\end{exercise}

\begin{exercise}[3]
Explain why full expansion can be computationally harmful.
\end{exercise}

\begin{exercise}[3]
Perform common subexpression elimination on an expression with repeated
factors.
\end{exercise}

\begin{exercise}[2]
Apply rewrite rules to simplify

\[
(x+0)(1)+0.
\]
\end{exercise}

\begin{exercise}[3]
Construct a pair of rewrite rules that creates a nonterminating loop.
\end{exercise}

\begin{exercise}[3]
Explain the meaning of confluence.
\end{exercise}

\begin{exercise}[2]
Simplify

\[
\sqrt{x^2}
\]

over the real numbers.
\end{exercise}

\begin{exercise}[3]
Explain how the assumption

\[
x>0
\]

changes the simplification.
\end{exercise}

\begin{exercise}[3]
Explain why

\[
\sqrt{ab}
=
\sqrt{a}\sqrt{b}
\]

requires care over complex numbers.
\end{exercise}

\begin{exercise}[3]
Write

\[
|x-2|
\]

as a piecewise function.
\end{exercise}

\begin{exercise}[3]
Solve

\[
x^2-4x+3\leq0
\]

symbolically.
\end{exercise}

\begin{exercise}[2]
Evaluate

\[
\sum_{k=1}^{n}k.
\]
\end{exercise}

\begin{exercise}[3]
Evaluate a finite geometric series.
\end{exercise}

\begin{exercise}[3]
Evaluate a telescoping sum.
\end{exercise}

\begin{exercise}[2]
Evaluate

\[
\prod_{k=1}^{n}k.
\]
\end{exercise}

\begin{exercise}[3]
Solve the recurrence

\[
a_{n+1}=3a_n.
\]
\end{exercise}

\begin{exercise}[3]
Solve a second-order constant-coefficient linear recurrence.
\end{exercise}

\begin{exercise}[2]
Solve

\[
y'=ay
\]

symbolically.
\end{exercise}

\begin{exercise}[3]
Solve a first-order separable ODE symbolically.
\end{exercise}

\begin{exercise}[3]
Derive the integrating-factor formula for a first-order linear ODE.
\end{exercise}

\begin{exercise}[3]
Compute a simple Laplace transform.
\end{exercise}

\begin{exercise}[3]
Use symbolic differentiation to construct a Jacobian matrix.
\end{exercise}

\begin{exercise}[3]
Construct the Hessian of a quadratic function.
\end{exercise}

\begin{exercise}[3]
Use symbolic derivatives to solve a one-variable optimization problem.
\end{exercise}

\begin{exercise}[3]
Verify a symbolic antiderivative by differentiation.
\end{exercise}

\begin{exercise}[3]
Verify a symbolic equation solution by substitution.
\end{exercise}

\begin{exercise}[3]
Explain symbolic-numeric computation.
\end{exercise}

\begin{exercise}[3]
Explain how exact preprocessing can improve a numerical computation.
\end{exercise}

\begin{exercise}[3]
Explain why expression-tree sharing can reduce memory use.
\end{exercise}

\begin{exercise}[3]
Explain lazy symbolic evaluation.
\end{exercise}

\begin{exercise}[3]
Distinguish arithmetic-operation complexity from bit complexity.
\end{exercise}

\begin{exercise}[4]
Study fast integer multiplication algorithms.
\end{exercise}

\begin{exercise}[4]
Study modular polynomial algorithms.
\end{exercise}

\begin{exercise}[4]
Study Chinese-remainder reconstruction in symbolic computation.
\end{exercise}

\begin{exercise}[4]
Study rational reconstruction.
\end{exercise}

\begin{exercise}[4]
Study fast polynomial multiplication.
\end{exercise}

\begin{exercise}[4]
Study subresultant polynomial remainder sequences.
\end{exercise}

\begin{exercise}[4]
Study square-free factorization algorithms.
\end{exercise}

\begin{exercise}[4]
Study polynomial factorization over finite fields.
\end{exercise}

\begin{exercise}[4]
Study factorization over the rational numbers.
\end{exercise}

\begin{exercise}[4]
Study algebraic-number arithmetic.
\end{exercise}

\begin{exercise}[4]
Study exact real root isolation.
\end{exercise}

\begin{exercise}[4]
Study resultant computation using Sylvester matrices.
\end{exercise}

\begin{exercise}[4]
Study Buchberger's algorithm for Gröbner bases.
\end{exercise}

\begin{exercise}[4]
Study monomial orderings and elimination theory.
\end{exercise}

\begin{exercise}[4]
Study symbolic integration algorithms for rational functions.
\end{exercise}

\begin{exercise}[4]
Study the Risch algorithm conceptually.
\end{exercise}

\begin{exercise}[4]
Study symbolic summation algorithms.
\end{exercise}

\begin{exercise}[4]
Study recurrence solving using generating functions.
\end{exercise}

\begin{exercise}[4]
Study symbolic asymptotic expansion algorithms.
\end{exercise}

\begin{exercise}[4]
Study simplification under complex branch conventions.
\end{exercise}

\begin{exercise}[4]
Study term rewriting, termination, and confluence.
\end{exercise}

\begin{exercise}[4]
Study symbolic expression DAGs and hash-consing.
\end{exercise}

\begin{exercise}[4]
Study common-subexpression elimination for automatic code generation.
\end{exercise}

\begin{exercise}[4]
Investigate symbolic-numeric equation solving.
\end{exercise}

%%%%%%%%%%%%%%%%%%%%%%%%%%%%%%%%%%%%%%%%%%%%%%%%%%%%%%%%%%%%
\subsection{Conceptual Summary}
\label{subsec:symbolic-computation-summary}
%%%%%%%%%%%%%%%%%%%%%%%%%%%%%%%%%%%%%%%%%%%%%%%%%%%%%%%%%%%%

Symbolic computation manipulates exact mathematical expressions rather
than immediately replacing them with numerical approximations.

A symbolic expression has structural form that can be represented using an
expression tree.

Symbolic methods include

\[
\boxed{
\text{simplification}
+
\text{expansion}
+
\text{factorization}
+
\text{exact arithmetic}.
}
\]

Polynomial computation includes:

\[
\boxed{
\text{division}
+
\text{gcd}
+
\text{factorization}
+
\text{elimination}.
}
\]

Calculus operations include:

\[
\boxed{
\text{differentiation}
+
\text{integration}
+
\text{limits}
+
\text{series}.
}
\]

Equation solving may combine:

\[
\boxed{
\text{substitution}
+
\text{resultants}
+
\text{polynomial methods}.
}
\]

Reliable symbolic manipulation requires attention to

\[
\boxed{
\text{domains}
+
\text{assumptions}
+
\text{branches}
+
\text{expression growth}.
}
\]

Symbolic and numerical methods are complementary:

\[
\boxed{
\text{symbolic structure}
+
\text{numerical evaluation}
=
\text{symbolic-numeric computation}.
}
\]

%%%%%%%%%%%%%%%%%%%%%%%%%%%%%%%%%%%%%%%%%%%%%%%%%%%%%%%%%%%%
\subsection{Section Connections}
\label{subsec:symbolic-computation-connections}
%%%%%%%%%%%%%%%%%%%%%%%%%%%%%%%%%%%%%%%%%%%%%%%%%%%%%%%%%%%%

\chapterconnection{
Symbolic Computation complements the approximate numerical methods of the
earlier sections of Chapter~11 by preserving exact mathematical
structure. It provides exact differentiation, algebraic simplification,
polynomial manipulation, equation solving, series expansion, transform
calculation, and symbolic preprocessing for numerical algorithms. These
methods connect strongly with algebra, number theory, calculus,
differential equations, optimization, and logic. The next section,
Computer Algebra, develops the broader algorithmic systems used to
implement exact algebraic computation, including polynomial arithmetic,
factorization, algebraic extensions, Gröbner bases, elimination, and the
architecture of computer algebra systems.
}

%%%%%%%%%%%%%%%%%%%%%%%%%%%%%%%%%%%%%%%%%%%%%%%%%%%%%%%%%%%%
% SECTION SOURCE TRACKING
%%%%%%%%%%%%%%%%%%%%%%%%%%%%%%%%%%%%%%%%%%%%%%%%%%%%%%%%%%%%

%------------------------------------------------------------
% 11.8 Symbolic Computation
%------------------------------------------------------------
%
% Source basis:
%   - Standard symbolic computation.
%   - Standard exact algebraic manipulation.
%   - Standard symbolic calculus and expression-rewriting concepts.
%
% Used for:
%   - symbolic versus numerical computation;
%   - expression structures and trees;
%   - syntactic and mathematical equality;
%   - canonical and normal forms;
%   - simplification;
%   - expansion;
%   - collection and factorization;
%   - exact integer and rational arithmetic;
%   - polynomial representations;
%   - polynomial arithmetic;
%   - polynomial division and gcd;
%   - square-free polynomials;
%   - repeated roots;
%   - polynomial factorization;
%   - rational-function simplification;
%   - partial fractions;
%   - symbolic differentiation;
%   - recursive differentiation;
%   - symbolic integration;
%   - substitution and integration by parts;
%   - nonelementary integrals;
%   - special functions;
%   - symbolic limits;
%   - Taylor, Laurent, and asymptotic series;
%   - symbolic equation solving;
%   - exact and numerical roots;
%   - radical-solvability limitations;
%   - root objects;
%   - systems of equations;
%   - substitution elimination;
%   - resultants;
%   - polynomial ideals and Gröbner-basis preview;
%   - symbolic matrices;
%   - determinants and inverses;
%   - symbolic eigenvalues;
%   - expression swell;
%   - common subexpression elimination;
%   - pattern matching;
%   - rewrite rules;
%   - termination and confluence;
%   - assumptions and domains;
%   - real and complex branches;
%   - piecewise expressions;
%   - symbolic inequalities;
%   - symbolic summation and products;
%   - recurrence relations;
%   - symbolic differential equations;
%   - symbolic Laplace and Fourier transforms;
%   - symbolic optimization;
%   - Jacobians and Hessians;
%   - code generation;
%   - symbolic-numeric computation;
%   - exact preprocessing;
%   - symbolic verification;
%   - proof assistance;
%   - computational and bit complexity;
%   - modular computation preview;
%   - evaluation-interpolation strategies;
%   - expression sharing;
%   - lazy evaluation;
%   - applications and exercises.
%
% Notes:
%   - Computer Algebra is developed next in Section 11.9.
%   - Mathematical Simulation follows in Section 11.10.
%   - High-Performance Mathematical Computing follows in Section 11.11.
%
%%%%%%%%%%%%%%%%%%%%%%%%%%%%%%%%%%%%%%%%%%%%%%%%%%%%%%%%%%%%